\UseRawInputEncoding
\documentclass[a4paper,12pt]{article}
\usepackage[margin=2cm]{geometry}
\usepackage[utf8]{inputenc}
\usepackage[T5]{fontenc}
\usepackage[vietnamese]{babel}
\usepackage{float}
\usepackage{graphicx}
\usepackage{titlesec}
\usepackage{setspace}
\usepackage{float}
\usepackage{tabularx}
\usepackage{longtable}
\usepackage[utf8]{inputenc}
\usepackage{listings}
\usepackage{listingsutf8}
\usepackage{xcolor}
\usepackage{mdframed}
\usepackage{listings}
\usepackage{caption}
% \usepackage{tcolorbox}
\usepackage[most]{tcolorbox}
\tcbuselibrary{breakable}

\usepackage{setspace, enumitem, titlesec}
\setstretch{0.95}
\setlength{\parskip}{0pt}
\setlist{nosep}
\titlespacing*{\section}{0pt}{0.5em}{0.3em}
\titlespacing*{\section}{0pt}{0.5em}{0.3em}
\titlespacing*{\subsection}{0pt}{0.3em}{0.2em}
\setlength{\jot}{2pt}
\setlist{itemsep=0pt, topsep=0pt, parsep=0pt, partopsep=0pt}
\setlength{\parskip}{0pt}      % bỏ khoảng cách giữa đoạn
\setlength{\parindent}{1em}    % tăng thụt đầu dòng cho dễ đọc


% ====== VS CODE LIGHT COLOR FIX ======
% ====== VS CODE LIGHT THEME ======
\definecolor{vscode-bg}{HTML}{EDEDED}   % nền code
\definecolor{vscode-gray}{HTML}{333333} % số dòng
\definecolor{vscode-blue}{HTML}{0066CC} % keyword
\definecolor{vscode-green}{HTML}{007700}% comment
\definecolor{vscode-red}{HTML}{B20000}  % string

\lstset{
  language=JavaScript,
  backgroundcolor=\color{vscode-bg},
  basicstyle=\ttfamily\small,
  keywordstyle=\color{vscode-blue}\bfseries,
  stringstyle=\color{vscode-red},
  commentstyle=\color{vscode-green}\itshape,
  numberstyle=\tiny\color{vscode-gray},
  numbers=left,
  numbersep=10pt,
  xleftmargin=20pt,
  frame=none,
  breaklines=true,
  columns=fullflexible,
  tabsize=2,
  showstringspaces=false
}

\begin{document}

% ================== TRANG BÌA ==================
\thispagestyle{empty}

% ======= TRƯỜNG - KHOA =======
\begin{center}
{\large \textbf{TRƯỜNG ĐẠI HỌC SÀI GÒN}}\\[2pt]
{\large \textbf{KHOA CÔNG NGHỆ THÔNG TIN}}\\[40pt]

% ======= LOGO =======
\includegraphics[width=4cm]{logoITSGU.png}\\[40pt]

% ======= TÊN MÔN =======
{\Large \textbf{KIỂM THỬ PHẦN MỀM}}\\[14pt]

\rule{12cm}{0.4pt}\\[24pt]


% ======= TÊN ĐỀ TÀI =======
{\normalsize Ứng dụng Đăng nhập \& Quản lý Sản phẩm}\\[4pt]
{\normalsize (Version 1.0)}\\[14pt]

\rule{12cm}{0.4pt}\\[24pt]

% ======= GVHD =======
\begin{tabular}{l l}
\textbf{GVHD:} & Từ Lăng Phiêu \\
\end{tabular}

% ======= SINH VIÊN THỰC HIỆN =======
\begin{center}
\textbf{Sinh viên thực hiện:}\\[4pt]
3123560086 - Dương Văn Thuật\\
3123530086 - Dương Văn A
\end{center}

% \textbf{GVHD:} \hspace{5pt} Từ Lăng Phiêu

\vfill

% ======= NGÀY THÁNG =======
{\normalsize TP. HỒ CHÍ MINH, THÁNG 11/2025}
\end{center}

\clearpage

% ================== MỤC LỤC ==================
\tableofcontents
\newpage

% ================== NỘI DUNG CHÍNH ==================

\section{Giới thiệu về Dự án}

\subsection{Tổng quan}

Dự án \textbf{FLoginFE\_BE} là một ứng dụng web hoàn chỉnh bao gồm:
\begin{itemize}
    \item \textbf{Chức năng Login}: Hệ thống đăng nhập với validation đầy đủ
    \item \textbf{Chức năng Product}: Quản lý sản phẩm (CRUD operations)
    \item \textbf{Frontend}: React 18+
    \item \textbf{Backend}: Spring Boot 3.2+
    \item \textbf{Testing}: Phát triển theo phương pháp TDD
\end{itemize}

\subsection{Công nghệ sử dụng}

\subsubsection{Frontend}
\begin{itemize}
    \item React 18+ -- Framework JavaScript
    \item React Testing Library -- Testing cho React
    \item Jest -- Testing framework
    \item Axios -- HTTP client
    \item CSS3 -- Styling với animations
\end{itemize}

\subsubsection{Backend}
\begin{itemize}
    \item Spring Boot 3.2+ -- Framework Java
    \item Java 17+
    \item JUnit 5 -- Testing framework
    \item Mockito -- Mock framework
    \item Maven -- Build tool
    \item Spring Data JPA -- Database operations
\end{itemize}

\subsection{Cấu trúc dự án}

\textbf{FLoginFE\_BE/}
\begin{itemize}
    \item \textbf{frontend/} - React Application
    \begin{itemize}
        \item{src/}
        \begin{itemize}
            \item components/ -- Login, Product components
            \item services/ -- API services
            \item utils/ -- Validation utilities
            \item tests/ -- Test files
        \end{itemize}
        \item{package.json}
    \end{itemize}

    \item \textbf{backend/} - Spring Boot API
    \begin{itemize}
        \item{src/}
        \begin{itemize}
            \item  main/java/com/flogin/
            \begin{itemize}
                \item controller/ - AuthController, ProductController
                \item service/ - Business logic
                \item dto/ - Data Transfer Objects
                \item entity/ - Database entities
                \item repository/ - Data access
            \end{itemize}
            \item test/java/ - Test files
        \end{itemize}
        \item{pom.xl}
    \end{itemize}
\end{itemize}

\section{Phân tích và Thiết kế Test Cases}

\subsection{Login - Phân tích và Test Scenarios}

\subsubsection{Yêu cầu}

Dựa trên chức năng đăng nhập (Login), hãy thực hiện:

a) Phân tích đầy đủ các yêu cầu chức năng của tính năng Login):
% =============================================================================================
\begin{table}[H]
\centering
\begin{tabularx}{\textwidth}{|>{\raggedright\arraybackslash}p{3.5cm}|>{\raggedright\arraybackslash}X|}

\hline
\textbf{Yêu cầu chức năng} & \textbf{Mô tả chi tiết} \\ \hline

1. Validation rules cho Username &
- Độ dài: 3–50 ký tự \newline
- Chỉ chứa: a–z, A–Z, 0–9, -, ., \_ \newline
- Không được rỗng \newline
- Không được có khoảng trắng đầu/cuối \newline
- Không được có ký tự đặc biệt ngoài danh sách cho phép \\ \hline

2. Validation rules cho Password &
- Độ dài: 6–100 ký tự \newline
- Phải chứa ít nhất 1 chữ số (0–9) \newline
- Phải chứa ít nhất 1 chữ cái \newline
- Có thể chứa ký tự đặc biệt (tùy hệ thống) \newline
- Không được rỗng \\ \hline

3. Authentication flow &
1. Người dùng nhập username + password \newline
2. Frontend validate $\rightarrow$ nếu lỗi $\rightarrow$ hiển thị thông báo \newline
3. Gửi request POST \texttt{/auth/login} \newline
4. Backend validate $\rightarrow$ trả về JWT token nếu hợp lệ \newline
4. Backend validate: \newline
    - Username có tồn tại không \newline
    - Password có khớp không (so sánh sau khi mã hóa) \newline
5. Nếu hợp lệ: \newline
    - Tạo JWT token hoặc session \newline
    - Trả về thông tin user (role, username, token) \newline
6. Nếu không hợp lệ $\rightarrow$ Trả về lỗi tương ứng \newline
7. Frontend nhận kết quả: \newline
    - Thành công $\rightarrow$ lưu token vào localStorage + redirect đến \texttt{/dashboard} \newline
    - Thất bại $\rightarrow$ hiển thị thông báo lỗi \\ \hline

4. Error handling &
- Username không hợp lệ: \newline
    + Nguyên nhân: Sai format <3 ký tự \newline
    + Phản hồi hệ thống: 400 Bad Request \newline
- Password không hợp lệ: \newline
    + Nguyên nhân: <6 ký tự hoặc không có chữ + số \newline
    + Phản hồi hệ thống: 400 Bad Request \newline
- User không tồn tại: \newline
    + Nguyên nhân: username không có trong Database \newline
    + Phản hồi hệ thống: 401 Unauthorized \newline
- Sai mật khẩu: \newline
    + Nguyên nhân: Password không khớp \newline
    + Phản hồi hệ thống: 401 Unauthorized \newline
- Tài khoản bị khóa: \newline
    + Nguyên nhân: Quá số lần login fail \newline
    + Phản hồi hệ thống: 403 Forbidden \newline
- Lỗi máy chủ: \newline
    + Nguyên nhân: Database lỗi, token lỗi \newline
    + Phản hồi hệ thống: 500 Internal Server Error \\ \hline

\end{tabularx}
\caption{Yêu cầu chức năng của tính năng đăng nhập}
\end{table}

% =============================================================================================

b) Liệt kê và mô tả ít nhất 10 test scenarios cho Login bao gồm (2 điểm):
% =============================================================================================
\begin{table}[H]
\centering
\begin{tabularx}{\textwidth}{|c|>{\raggedright\arraybackslash}X|c|>{\raggedright\arraybackslash}X|>{\raggedright\arraybackslash}X|}
\hline
\textbf{ID} & \textbf{Test Scenario} & \textbf{Type} & \textbf{Input} & \textbf{Expected Result} \\ \hline

LOGIN\_TC01 & Đăng nhập thành công & Happy path &
- Username: admin \newline
- Password: admin123  & Login thành công, trả về token chuyển trang. \\
\hline
LOGIN\_TC02 & Username rỗng & Negative tests &
- Username: "" \newline
- Password hợp lệ  & Hiển thị lỗi "Username không được để trống" \\
\hline
LOGIN\_TC03 & Password rỗng & Negative tests &
- Username hợp lệ \newline
- Password: ""  & Hiển thị lỗi "Password không được để trống \\
\hline
LOGIN\_TC04 & Username không đúng format & Negative tests &
- Username: abc@123 \newline
- Password hợp lệ  & Lỗi "Username không hợp lệ" \\
\hline
LOGIN\_TC05 & Đăng nhập thành công & Negative tests &
- Username không tồn tại \newline
- Password hợp lệ  & Lỗi "User không tồn tại" \\
\hline
LOGIN\_TC06 & Sai password & Negative tests &
- Username tồn tại \newline
- Password sai & Lỗi "Sai mật khẩu" \\
\hline
LOGIN\_TC07 & Username < 3 ký tự & Boundary tests &
- Username: ab \newline
- Password hợp lệ  & Lỗi "Username quá ngắn" \\
\hline
LOGIN\_TC08 & Password < 6 ký tự & Boundary tests &
- Username hợp lệ \newline
- Password: abc12  & Lỗi "Password quá ngắn" \\
\hline
LOGIN\_TC09 & Username 50 ký tự (max) & Boundary tests &
- Username: Username = 50 ký tự hợp lệ \newline
- Password hợp lệ  & Chấp nhận $\rightarrow$ login bình thường (nếu user tồn tại) \\
\hline
LOGIN\_TC10 & Username có khoảng trắng đầu/cuối & Edge cases &
- Username: "admin " \newline
- Password hợp lệ  & Lỗi "Username không hợp lệ" \\
\hline
LOGIN\_TC11 & Password hợp lệ nhưng chứa nhiều ký tự đặc biệt & Edge cases &
- Username hợp lệ \newline
- Password: "abc123!!!@@@"  & Login OK (miễn có chữ + số + đủ độ dài) \\
\hline
LOGIN\_TC12 & Username với ký tự gạch dưới/dấu chấm & Edge cases &
- Username: "user.test\_01" \newline
- Password hợp lệ  & Cho phép (valid format) \\
\hline
\end{tabularx}
\caption{Danh sách 12 Test Scenarios cho chức năng Login}
\end{table}
% =============================================================================================
c) Phân loại test scenarios theo mức độ ưu tiên (Critical, High, Medium, Low) và giải thích (1 điểm)
% =============================================================================================
\begin{table}[h!]
\centering
\begin{tabularx}{\textwidth}{|p{2.2cm}|>{\raggedright\arraybackslash}X|>{\raggedright\arraybackslash}X|}
\hline
\textbf{Priority} & \textbf{Test Scenario} & \textbf{Lý do} \\
\hline

\textbf{Critical} &
TS01 – Đăng nhập thành công \newline
TS05 – Sai username \newline
TS06 – Sai password &
Ảnh hưởng trực tiếp đến luồng chính và bảo mật. Nếu sai sẽ không thể login hoặc gây lỗ hổng truy cập trái phép. \\
\hline

\textbf{High} &
TS02 – Username rỗng \newline
TS03 – Password rỗng \newline
TS08 – Password < 6 ký tự \newline
TS11 – Password chứa nhiều ký tự đặc biệt \newline
TS12 – Username có dấu chấm/gạch dưới &
Validation quan trọng, liên quan đến bảo mật (password yếu, rỗng) hoặc format phổ biến cần hỗ trợ đúng. \\
\hline

\textbf{Medium} &
TS04 – Username sai format (ký tự lạ) \newline
TS07 – Username < 3 ký tự \newline
TS10 – Username có khoảng trắng đầu/cuối &
Các trường hợp validation/edge case, ít ảnh hưởng trực tiếp đến bảo mật nhưng cần kiểm soát để tránh lỗi nhập liệu. \\
\hline

\textbf{Low} &
TS09 – Username 50 ký tự (max) &
Trường hợp hiếm gặp, kiểm tra giới hạn hệ thống, ít ảnh hưởng đến UX và bảo mật. \\
\hline
\end{tabularx}
\caption{Phân loại Test Scenarios (TS01–TS12) theo mức độ ưu tiên}
\end{table}
% =============================================================================================
\subsubsection{Thiết kế 5 test cases quan trọng nhất cho Login theo template: }
% =============================================================================================
\begin{table}[H]
\centering
\begin{tabularx}{\textwidth}{|p{3.2cm}|>{\raggedright\arraybackslash}X|}
\hline
\textbf{Thuộc tính} & \textbf{Chi tiết} \\ \hline

\textbf{Test Case ID} & TC\_LOGIN\_001 \\ \hline

\textbf{Tên Test} & Đăng nhập thành công với credentials hợp lệ \\ \hline

\textbf{Mức độ ưu tiên} & Critical \\ \hline

\textbf{Điều kiện tiên quyết} &
- User account exists \newline
- Application is running \\ \hline

\textbf{Các bước kiểm thử} &
1. Navigate to login page \newline
2. Enter valid username \newline
3. Enter valid password \newline
4. Click Login button \\ \hline

\textbf{Dữ liệu kiểm thử} &
- Username: testuser \newline
- Password: Test123 \\ \hline

\textbf{Kết quả mong đợi} &
- Success message displayed \newline
- Token stored \newline
- Redirect to dashboard \\ \hline

\textbf{Kết quả thực tế} &
(Để trống) \\ \hline

\textbf{Trạng thái} &
Not Run \\ \hline

\end{tabularx}
\caption{Test Case TC\_LOGIN\_001 – Đăng nhập thành công với credentials hợp lệ}
\end{table}
% =============================================================================================
\begin{table}[H]
\centering
\begin{tabularx}{\textwidth}{|p{3.2cm}|>{\raggedright\arraybackslash}X|}
\hline
\textbf{Thuộc tính} & \textbf{Chi tiết} \\ \hline

\textbf{Test Case ID} & TC\_LOGIN\_002 \\ \hline

\textbf{Tên Test} & Đăng nhập thất bại khi username trống \\ \hline

\textbf{Mức độ ưu tiên} & High \\ \hline

\textbf{Điều kiện tiên quyết} &
- Ứng dụng đang chạy\\ \hline

\textbf{Các bước kiểm thử} &
1. Truy cập trang /login \newline
2. Để trống trường username \newline
3. Nhập password hợp lệ \newline
4. Click nút Login \\ \hline

\textbf{Dữ liệu kiểm thử} &
Username: (trống) \newline
Password: Test123 \\ \hline

\textbf{Kết quả mong đợi} &
- Hiển thị thông báo: "Username không được để trống" \newline
- Không gọi API Login \\ \hline

\textbf{Kết quả thực tế} &
(Để trống) \\ \hline

\textbf{Trạng thái} &
Not Run \\ \hline

\end{tabularx}
\caption{Test Case TC\_LOGIN\_002 – Đăng nhập thất bại khi username trống}
\end{table}
% =============================================================================================
\begin{table}[H]
\centering
\begin{tabularx}{\textwidth}{|p{3.2cm}|>{\raggedright\arraybackslash}X|}
\hline
\textbf{Thuộc tính} & \textbf{Chi tiết} \\ \hline

\textbf{Test Case ID} & TC\_LOGIN\_003 \\ \hline

\textbf{Tên Test} & Đăng nhập thất bại khi password trống \\ \hline

\textbf{Mức độ ưu tiên} & High \\ \hline

\textbf{Điều kiện tiên quyết} &
- Ứng dụng đang chạy \\ \hline

\textbf{Các bước kiểm thử} &
1. Truy cập trang /login \newline
2. Nhập username hợp lệ \newline
3. Để trống trường password \newline
4. Click nút Login \\ \hline

\textbf{Dữ liệu kiểm thử} &
Username: testuser \newline
Password: (trống) \\ \hline

\textbf{Kết quả mong đợi} &
- Hiển thị thông báo: "Password is required" \newline
- Không gửi request login \\ \hline

\textbf{Kết quả thực tế} &
(Để trống) \\ \hline

\textbf{Trạng thái} &
Not Run \\ \hline

\end{tabularx}
\caption{Test Case TC\_LOGIN\_003 – Đăng nhập thất bại khi password trống}
\end{table}
% =============================================================================================
\begin{table}[H]
\centering
\begin{tabularx}{\textwidth}{|p{3.2cm}|>{\raggedright\arraybackslash}X|}
\hline
\textbf{Thuộc tính} & \textbf{Chi tiết} \\ \hline

\textbf{Test Case ID} & TC\_LOGIN\_004 \\ \hline

\textbf{Tên Test} & Đăng nhập thất bại khi mật khẩu không đúng \\ \hline

\textbf{Mức độ ưu tiên} & Critical \\ \hline

\textbf{Điều kiện tiên quyết} &
- User account exists \newline
- Application is running \\ \hline

\textbf{Các bước kiểm thử} &
1. Navigate to login page \newline
2. Enter valid username \newline
3. Enter wrong password \newline
4. Click Login button \\ \hline

\textbf{Dữ liệu kiểm thử} &
Username: testuser \newline
Password: Wrong Pass \\ \hline

\textbf{Kết quả mong đợi} &
- Message: "Invalid username or password" \newline
- Không đăng nhập thành công \newline
- Không tạo token \\ \hline

\textbf{Kết quả thực tế} &
(Để trống) \\ \hline

\textbf{Trạng thái} &
Not Run \\ \hline

\end{tabularx}
\caption{Test Case TC\_LOGIN\_004 – Đăng nhập thất bại khi mật khẩu không đúng}
\end{table}
% =============================================================================================
\begin{table}[H]
\centering
\begin{tabularx}{\textwidth}{|p{3.2cm}|>{\raggedright\arraybackslash}X|}
\hline
\textbf{Thuộc tính} & \textbf{Chi tiết} \\ \hline

\textbf{Test Case ID} & TC\_LOGIN\_005 \\ \hline

\textbf{Tên Test} & Đăng nhập thất bại khi username sai định dạng \\ \hline

\textbf{Mức độ ưu tiên} & High \\ \hline

\textbf{Điều kiện tiên quyết} &
- Ứng dụng đang chạy \\ \hline

\textbf{Các bước kiểm thử} &
1. Navigate to login page \newline
2. Enter username có ký tự không hợp lệ \newline
3. Enter valid password \newline
4. Click Login button \\ \hline

\textbf{Dữ liệu kiểm thử} &
Username: test!@@@ \newline
Password: Test123 \\ \hline

\textbf{Kết quả mong đợi} &
- Message: "Invalid username format" \newline
- Không gửi request login \\ \hline

\textbf{Kết quả thực tế} &
(Để trống) \\ \hline

\textbf{Trạng thái} &
Not Run \\ \hline

\end{tabularx}
\caption{Test Case TC\_LOGIN\_005 – Đăng nhập thất bại khi username sai định dạng}
\end{table}
% =============================================================================================

\subsection{Product – Phân tích và Test Scenarios}

\subsubsection{Yêu cầu}

Dựa trên chức năng quản lý sản phẩm (Product Management), hãy thực hiện:

a) Phân tích đầy đủ các yêu cầu chức năng của Product CRUD:
% ============================================================================================
\begin{table}[H]
\centering
\begin{tabular}{|p{3.5cm}|p{10.2cm}|}
\hline
\textbf{Chức năng} & \textbf{Yêu cầu chức năng chi tiết} \\ \hline

\textbf{Create} &
- Người dùng (đã đăng nhập) truy cập form thêm sản phẩm \newline
- Nhập các trường: \texttt{Nam}, \texttt{Price}, \texttt{Quantity}, \texttt{Description}, \texttt{Category} (tùy chọn) \newline
- Validation: \newline
    + Name: 3-100 ký tự, không rỗng \newline
    + Price: >0, $\leq$999,999,999 \newline
    + Quantity: 0–99,999 \newline
    + Description: $\leq$500 ký tự \newline
    + Category: phải thuộc danh sách có sẵn \newline
- Success: sản phẩm được thêm vào danh sách \\ \hline

\textbf{Read} &
- Hiển thị danh sách sản phẩm \newline
- Hiển thị chi tiết sản phẩm \\ \hline

\textbf{Update} &
- Chọn sản phẩm, chỉnh sửa thông tin \newline
- Validation tương tự Create\\ \hline

\textbf{Delete} &
- Xóa sản phẩm đã tồn tại \newline
- Xử lý xóa sản phẩm không tồn tại hoặc đang được tham chiếu\\ \hline

\end{tabular}
\caption{ Chức năng CRUD và yêu cầu chức năng chi tiết}
\end{table}
% =============================================================================================
b) Liệt kê và mô tả ít nhất 10 test scenarios cho Product, bao gồm:
% =============================================================================================
\begin{table}[H]
\centering
\begin{tabularx}{\textwidth}{|c|>{\raggedright\arraybackslash}X|c|>{\raggedright\arraybackslash}X|>{\raggedright\arraybackslash}X|}
\hline
\textbf{ID} & \textbf{Test Scenario} & \textbf{Type} & \textbf{Input} & \textbf{Expected Result} \\ \hline

CRUD\_TC01 & Thêm sản phẩm thành công & Happy path &
- Name: "Bút bi Thiên Long 0.5mm" \newline
- Price: 5000 \newline
- Quantity: 500 \newline
- Category: "Books" & Sản phẩm thêm thành công, hiển thị trong danh sách \\ \hline

CRUD\_TC02 & Cập nhật sản phẩm thành công & Happy path &
- Chọn sản phẩm đã tồn tại \newline
- Thay đổi Name, Price, Quantity hợp lệ & Sản phẩm được cập nhật thành công, hiển thị thông báo \\ \hline

CRUD\_TC03 & Thêm sản phẩm với Name rỗng & Negative tests &
- Name: "" \newline
- Price: 500000 \newline
- Quantity: 10 \newline
- Category: "Books" & Lỗi "Tên sản phẩm không được để trống" \\ \hline

CRUD\_TC04 & Thêm sản phẩm với Price $\leq 0$ & Negative tests &
- Name: "Vở ô ly 100 trang" \newline
- Price: 0 \newline
- Quantity: 10 \newline
- Category: "Books" & Lỗi "Giá phải lớn hơn 0" \\ \hline

CRUD\_TC05 & Thêm sản phẩm với Quantity âm & Negative tests &
- Name: "Máy tính Casio FX-580VN X" \newline
- Price: 500000 \newline
- Quantity: -5 \newline
- Category: "Books" & Lỗi "Số lượng không được âm" \\ \hline

CRUD\_TC06 & Name = 3 ký tự (min) & Boundary tests &
- Name: "Bút" \newline
- Price: 500000 \newline
- Quantity: 10 \newline
- Category: "Books" & Hợp lệ, sản phẩm thêm thành công \\ \hline

CRUD\_TC07 & Name = 100 ký tự (max) & Boundary tests &
- Name: "aaaaaaaa...a" (100 ký tự) \newline
- Price: 500000 \newline
- Quantity: 10 \newline
- Category: "Books" & Hợp lệ, sản phẩm thêm thành công \\ \hline

\end{tabularx}
\end{table}

\begin{table}[H]
\centering
\begin{tabularx}{\textwidth}{|c|>{\raggedright\arraybackslash}X|c|>{\raggedright\arraybackslash}X|>{\raggedright\arraybackslash}X|}
\hline
\textbf{ID} & \textbf{Test Scenario} & \textbf{Type} & \textbf{Input} & \textbf{Expected Result} \\ \hline

CRUD\_TC08 & Price = 999999999 (max) & Boundary tests &
- Name: "Tai nghe Bluetooth Sony WH-CH520" \newline
- Price: 999999999 \newline
- Quantity: 10 \newline
- Category: "Electronics" & Hợp lệ, sản phẩm thêm thành công \\ \hline

CRUD\_TC09 & Thêm sản phẩm trùng tên & Edge cases &
- Name: "USB 3.0 64GB Kingston" \newline
- Price: 199000 \newline
- Quantity: 50 \newline
- Category: "Electronics" & Lỗi "Sản phẩm đã tồn tại" hoặc chặn trùng \\ \hline

CRUD\_TC10 & Thêm sản phẩm Name chứa khoảng trắng đầu/cuối & Edge cases &
- Name: " Áo thun nam Polo cơ bản" \newline
- Price: 25000000 \newline
- Quantity: 50 \newline
- Category: "Clothing" & Tự trim khoảng trắng hoặc hiển thị lỗi "Tên không hợp lệ" \\ \hline

CRUD\_TC11 & Xóa sản phẩm thành công & Happy path &
- Chọn sản phẩm đã tồn tại \newline
- Click nút Delete & Sản phẩm bị xóa khỏi danh sách, hiển thị thông báo thành công \\ \hline

CRUD\_TC12 & Xóa sản phẩm không tồn tại & Edge cases &
- Nhập ID sản phẩm không tồn tại \newline
- Gọi API Delete & Hiển thị lỗi "Sản phẩm không tồn tại", không ảnh hưởng dữ liệu hiện có \\ \hline

\end{tabularx}
\caption{Danh sách 12 Test Scenarios cho chức năng CRUD sản phẩm}
\end{table}
% =============================================================================================
c) Phân loại test scenarios theo mức độ ưu tiên và giải thích.
% ===========================================================================================
\textbf{Validation Rules cho Product}

% ===========================================================================================
\begin{table}[H]
\centering
\begin{tabularx}{\textwidth}{|p{2.5cm}|>{\raggedright\arraybackslash}X|>{\raggedright\arraybackslash}X|}
\hline
\textbf{Priority} & \textbf{Test Scenarios} & \textbf{Giải thích} \\ \hline

\textbf{Critical} &
- TS01: Thêm sản phẩm thành công \newline
- TS02: Cập nhật sản phẩm thành công \newline
- TS11: Xóa sản phẩm thành công &
Các luồng chính của CRUD. Nếu không hoạt động thì hệ thống không usable. \\ \hline

\textbf{High} &
- TS03: Thêm sản phẩm với Name rỗng \newline
- TS04: Thêm sản phẩm với Price $\leq 0$ \newline
- TS05: Thêm sản phẩm với Quantity âm \newline
- TS12: Xóa sản phẩm không tồn tại &
Validation quan trọng, ảnh hưởng trực tiếp đến tính đúng đắn và bảo mật dữ liệu. \\ \hline

\textbf{Medium} &
- TS06: Name = 3 ký tự (min) \newline
- TS07: Name = 100 ký tự (max) \newline
- TS08: Price = 999999999 (max) &
Boundary tests kiểm tra giới hạn hệ thống. Quan trọng nhưng ít xảy ra trong thực tế. \\ \hline

\textbf{Low} &
- TS09: Thêm sản phẩm trùng tên \newline
- TS10: Thêm sản phẩm Name chứa khoảng trắng đầu/cuối &
Edge cases hiếm gặp, ảnh hưởng chủ yếu đến UX hoặc tính toàn vẹn dữ liệu nhưng không gây lỗi nghiêm trọng. \\ \hline

\end{tabularx}
\caption{Phân loại 12 Test Scenarios CRUD sản phẩm theo mức độ ưu tiên}
\end{table}

% ===========================================================================================

\subsubsection{Thiết kế Test Cases chi tiết (5 điểm)}
Thiết kế 5 test cases quan trọng nhất cho Product Management:
% ===========================================================================================

\begin{table}[H]
\centering
\begin{tabularx}{\textwidth}{|p{3cm}|>{\raggedright\arraybackslash}X|}
\hline
\textbf{Thuộc tính} & \textbf{Chi tiết} \\ \hline

\textbf{Test Case ID} & TC\_PRODUCT\_001 \\ \hline
\textbf{Tên Test} & Tạo sản phẩm mới thành công \\ \hline
\textbf{Mức độ ưu tiên} & Critical \\ \hline
\textbf{Điều kiện tiên quyết} &
- User đã đăng nhập \newline
- User có quyền tạo sản phẩm \\ \hline
\textbf{Các bước kiểm thử} &
1. Navigate to Product page \newline
2. Click "Add New Product" \newline
3. Enter product information \newline
4. Click Save button \\ \hline
\textbf{Dữ liệu kiểm thử} &
Name: Laptop Dell \newline
Price: 15000000 \newline
Quantity: 10 \newline
Category: Electronics \\ \hline
\textbf{Kết quả mong đợi} &
- Product created successfully \newline
- Success message displayed \newline
- Product appears in list \\ \hline
\textbf{Kết quả thực tế} & (Để trống) \\ \hline
\textbf{Trạng thái} & Not Run \\ \hline

\end{tabularx}
\caption{Test Case TC\_PRODUCT\_001 – Tạo sản phẩm mới thành công}
\end{table}

% ===========================================================================================

\begin{table}[H]
\centering
\begin{tabularx}{\textwidth}{|p{3cm}|>{\raggedright\arraybackslash}X|}
\hline
\textbf{Thuộc tính} & \textbf{Chi tiết} \\ \hline

\textbf{Test Case ID} & TC\_PRODUCT\_002 \\ \hline
\textbf{Tên Test} & Cập nhật sản phẩm thành công \\ \hline
\textbf{Mức độ ưu tiên} & Critical \\ \hline
\textbf{Điều kiện tiên quyết} &
- User đã đăng nhập \newline
- User có quyền cập nhật sản phẩm \newline
- Sản phẩm "Laptop Dell" đã tồn tại trong danh sách \\ \hline
\textbf{Các bước kiểm thử} &
1. Navigate to Product page \newline
2. Select product "Laptop Dell" \newline
3. Click "Edit" button \newline
4. Update product information (ví dụ: Price, Quantity) \newline
5. Click Save button \\ \hline
\textbf{Dữ liệu kiểm thử} &
Name: Laptop Dell \newline
Price: 16000000 \newline
Quantity: 15 \newline
Category: Electronics \\ \hline
\textbf{Kết quả mong đợi} &
- Product updated successfully \newline
- Success message displayed \newline
- Product list reflects updated values \\ \hline
\textbf{Kết quả thực tế} & (Để trống) \\ \hline
\textbf{Trạng thái} & Not Run \\ \hline

\end{tabularx}
\caption{Test Case TC\_PRODUCT\_002 – Cập nhật sản phẩm thành công}
\end{table}

% ===========================================================================================

\begin{table}[H]
\centering
\begin{tabularx}{\textwidth}{|p{3cm}|>{\raggedright\arraybackslash}X|}
\hline
\textbf{Thuộc tính} & \textbf{Chi tiết} \\ \hline

\textbf{Test Case ID} & TC\_PRODUCT\_003 \\ \hline
\textbf{Tên Test} & Xóa sản phẩm thành công \\ \hline
\textbf{Mức độ ưu tiên} & Critical \\ \hline
\textbf{Điều kiện tiên quyết} &
- User đã đăng nhập \newline
- User có quyền xóa sản phẩm \newline
- Sản phẩm "Laptop Dell" đã tồn tại \\ \hline
\textbf{Các bước kiểm thử} &
1. Navigate to Product page \newline
2. Select product "Laptop Dell" \newline
3. Click "Delete" button \newline
4. Confirm deletion \\ \hline
\textbf{Dữ liệu kiểm thử} &
Product: Laptop Dell \\ \hline
\textbf{Kết quả mong đợi} &
- Product deleted successfully \newline
- Success message displayed \newline
- Product không còn xuất hiện trong danh sách \\ \hline
\textbf{Kết quả thực tế} & (Để trống) \\ \hline
\textbf{Trạng thái} & Not Run \\ \hline

\end{tabularx}
\caption{Test Case TC\_PRODUCT\_003 – Xóa sản phẩm thành công}
\end{table}

% ===========================================================================================
\begin{table}[H]
\centering
\begin{tabularx}{\textwidth}{|p{3cm}|>{\raggedright\arraybackslash}X|}
\hline
\textbf{Thuộc tính} & \textbf{Chi tiết} \\ \hline

\textbf{Test Case ID} & TC\_PRODUCT\_004 \\ \hline
\textbf{Tên Test} & Thêm sản phẩm với Name rỗng \\ \hline
\textbf{Mức độ ưu tiên} & High \\ \hline
\textbf{Điều kiện tiên quyết} &
- User đã đăng nhập \newline
- User có quyền tạo sản phẩm \\ \hline
\textbf{Các bước kiểm thử} &
1. Navigate to Product page \newline
2. Click "Add New Product" \newline
3. Để trống trường Name \newline
4. Nhập các thông tin khác hợp lệ \newline
5. Click Save button \\ \hline
\textbf{Dữ liệu kiểm thử} &
Name: "" \newline
Price: 500000 \newline
Quantity: 10 \newline
Category: Books \\ \hline
\textbf{Kết quả mong đợi} &
- Hiển thị lỗi "Tên sản phẩm không được để trống" \newline
- Không gọi API tạo sản phẩm \\ \hline
\textbf{Kết quả thực tế} & (Để trống) \\ \hline
\textbf{Trạng thái} & Not Run \\ \hline

\end{tabularx}
\caption{Test Case TC\_PRODUCT\_004 – Thêm sản phẩm với Name rỗng}
\end{table}

% ===========================================================================================
\begin{table}[H]
\centering
\begin{tabularx}{\textwidth}{|p{3cm}|>{\raggedright\arraybackslash}X|}
\hline
\textbf{Thuộc tính} & \textbf{Chi tiết} \\ \hline

\textbf{Test Case ID} & TC\_PRODUCT\_005 \\ \hline
\textbf{Tên Test} & Thêm sản phẩm với Price $\leq 0$ \\ \hline
\textbf{Mức độ ưu tiên} & High \\ \hline
\textbf{Điều kiện tiên quyết} &
- User đã đăng nhập \newline
- User có quyền tạo sản phẩm \\ \hline
\textbf{Các bước kiểm thử} &
1. Navigate to Product page \newline
2. Click "Add New Product" \newline
3. Nhập Name hợp lệ \newline
4. Nhập Price = 0 hoặc âm \newline
5. Nhập các thông tin khác hợp lệ \newline
6. Click Save button \\ \hline
\textbf{Dữ liệu kiểm thử} &
Name: "Vở ô ly 100 trang" \newline
Price: 0 \newline
Quantity: 10 \newline
Category: Books \\ \hline
\textbf{Kết quả mong đợi} &
- Hiển thị lỗi "Giá phải lớn hơn 0" \newline
- Không gọi API tạo sản phẩm \\ \hline
\textbf{Kết quả thực tế} & (Để trống) \\ \hline
\textbf{Trạng thái} & Not Run \\ \hline

\end{tabularx}
\caption{Test Case TC\_PRODUCT\_005 – Thêm sản phẩm với Price $\leq 0$}
\end{table}

% ===========================================================================================

\section{Unit Testing và Test-Driven Development}

\subsection{Login - Unit Tests Frontend và Backend}

\subsubsection{ Frontend Unit Tests - Validation Login}

Áp dụng TDD để phát triển unit tests cho validation module của Login.

\textbf{Yêu cầu:}

a) Viết unit tests cho \texttt{validateUsername()} (2 điểm):

\begin{tcolorbox}[breakable, colback=green!5,colframe=black!50,title=validation.test.js]
\begin{lstlisting}
import { validateUsername, validatePassword } from './validation';

describe('Login Validation Tests', () => {
    describe('validateUsername()', () => {
        // Test username rong
        test('TC1: Username rong nen tra ve loi', () => {
            // Day la loi mong muon neu ham tra ve chuoi loi, hoac true/false
            expect(validateUsername('')).toBe('Ten dang nhap khong duoc de trong');
        });
        // Test username qua ngan (Boundary Test: < 3 ky tu)
        test('TC2: Username qua ngan (2 ky tu) nen tra ve loi', () => {
            expect(validateUsername('ab')).toBe('Ten dang nhap phai co it nhat 3 ky tu');
        });
        // Test username qua dai (Boundary Test: > 50 ky tu)
        test('TC3: Username qua dai (51 ky tu) nen tra ve loi', () => {
            // Chuoi 51 ky tu
            const longUsername = 'a'.repeat(51);
            expect(validateUsername(longUsername)).toBe('Ten dang nhap khong duoc vuot qua 50 ky tu');
        });
        // Test ky tu dac biet khong hop le
        test('TC4: Username chua ky tu dac biet (@, #) nen tra ve loi', () => {
            expect(validateUsername('user@123')).toBe('Ten dang nhap chi duoc chua a-z, A-Z, 0-9');
        });
        // Test username hop le
        test('TC5: Username hop le (user123) nen khong co loi', () => {
            expect(validateUsername('user123')).toBe('');
        });
    });
\end{lstlisting}
\end{tcolorbox}

b) Unit Tests cho validatePassword() (2 diem)
\begin{tcolorbox}[breakable, colback=green!5,colframe=black!50,title=validation.test.js]
\begin{lstlisting}

    describe('validatePassword()', () => {
        // Test password rong
        test('TC6: Password rong nen tra ve loi', () => {
            expect(validatePassword('')).toBe('Mat khau khong duoc de trong');
        });

        // Test password qua ngan (Boundary Test: < 6 ky tu)
        test('TC7: Password qua ngan (5 ky tu) nen tra ve loi', () => {
            expect(validatePassword('Pass1')).toBe('Mat khau phai co it nhat 6 ky tu');
        });

        // Test password khong co chu hoac so
        test('TC8: Password khong co so nen tra ve loi', () => {
            expect(validatePassword('testtest')).toBe('Mat khau phai chua ca chu va so');
        });

        test('TC9: Password khong co chu nen tra ve loi', () => {
            expect(validatePassword('1234567')).toBe('Mat khau phai chua ca chu va so');
        });

        // Test password hop le
        test('TC10: Password hop le (Test123) nen khong co loi', () => {
            expect(validatePassword('Test123')).toBe('');
        });
    });
});

\end{lstlisting}
\end{tcolorbox}


c) Coverage $\geq 90\%$ cho validation module



\subsubsection{Backend Unit Tests - Login Service}

Viết unit tests cho \texttt{AuthService} của backend:

\textbf{Yêu cầu:}

a) Test method \texttt{authenticate()} với các scenarios :
\begin{tcolorbox}[breakable, colback=green!5,colframe=black!50,title=AuthServiceTest.java]
\begin{lstlisting}

    package com.flogin.service.impl;

import com.flogin.dto.LoginRequest;
import com.flogin.dto.LoginResponse;
import com.flogin.service.JwtService;

import org.junit.jupiter.api.DisplayName;
import org.junit.jupiter.api.Test;
import org.junit.jupiter.api.extension.ExtendWith;
import org.mockito.InjectMocks;
import org.mockito.Mock;
import org.mockito.junit.jupiter.MockitoExtension;
import org.springframework.security.authentication.AuthenticationManager;
import org.springframework.security.authentication.UsernamePasswordAuthenticationToken;
import org.springframework.security.core.userdetails.UsernameNotFoundException;
import org.springframework.security.authentication.BadCredentialsException;

import static org.junit.jupiter.api.Assertions.*;
import static org.mockito.Mockito.*;

@ExtendWith(MockitoExtension.class)
@DisplayName("Cau 2.1.2: AuthService Unit Tests")
class AuthServiceTest {

    @Mock
    private AuthenticationManager authManager;

    @Mock
    private JwtService jwtService;

    @InjectMocks
    private AuthServiceImpl authService;

    // 1. Login thanh cong
    @Test
    @DisplayName("TC1: Login thanh cong va tao JWT Token")
    void testLoginSuccess() {
        LoginRequest request = new LoginRequest("testuser", "Test123");

        when(jwtService.generateToken("testuser")).thenReturn("mock-token");

        LoginResponse response = authService.authenticate(request);

        assertTrue(response.isSuccess());
        assertEquals("Dang nhap thanh cong", response.getMessage());
        assertEquals("mock-token", response.getToken());
\end{lstlisting}
\end{tcolorbox}
\begin{tcolorbox}[breakable, colback=green!5,colframe=black!50,title=AuthServiceTest.java]
\begin{lstlisting}
        verify(authManager, times(1))
                .authenticate(any(UsernamePasswordAuthenticationToken.class));
    }

    // 2. Username khong ton tai
    @Test
    @DisplayName("TC2: Login that bai do username khong ton tai")
    void testLoginFailureUserNotFound() {
        LoginRequest request = new LoginRequest("nouser", "123456");

        doThrow(new UsernameNotFoundException("not found"))
                .when(authManager)
                .authenticate(any(UsernamePasswordAuthenticationToken.class));

        LoginResponse response = authService.authenticate(request);

        assertFalse(response.isSuccess());
        assertEquals("Sai thong tin dang nhap", response.getMessage());
        assertNull(response.getToken());

        verify(jwtService, never()).generateToken(any());
    }

    // 3. Password sai
    @Test
    @DisplayName("TC3: Login that bai do password sai")
    void testWrongPassword() {
        LoginRequest request = new LoginRequest("abc", "wrongpass");

        doThrow(new BadCredentialsException("bad credentials"))
                .when(authManager)
                .authenticate(any(UsernamePasswordAuthenticationToken.class));

        LoginResponse response = authService.authenticate(request);

        assertFalse(response.isSuccess());
        assertNull(response.getToken());
        assertEquals("Sai thong tin dang nhap", response.getMessage());
    }
\end{lstlisting}
\end{tcolorbox}
\begin{tcolorbox}[breakable, colback=green!5,colframe=black!50,title=AuthServiceTest.java]
\begin{lstlisting}
    // 4. Validation errors
    @Test
    @DisplayName("TC4: Validation loi - username null")
    void testValidationError() {
        LoginRequest request = new LoginRequest(null, "123");

        LoginResponse response = authService.authenticate(request);

        assertFalse(response.isSuccess());
        assertEquals("Thieu username", response.getMessage());
        assertNull(response.getToken());

        verify(authManager, never()).authenticate(any());
    }
}
\end{lstlisting}
\end{tcolorbox}

b) Test validation methods riêng lẻ (1 điểm)

c) Coverage $\geq 85\%$ cho AuthService (1 điểm)


% ===============================================================================================================================
\subsection{Product - Unit Tests Frontend và Backend}

\subsubsection{Frontend Unit Tests - Product Validation}

Áp dụng TDD cho validation của Product:

\textbf{Yêu cầu:}

a) Viết unit tests cho \texttt{validateProduct()}:

\begin{tcolorbox}[breakable, colback=green!5,colframe=black!50,title=productValidation.test.js]
\begin{lstlisting}
    import { validateProduct } from "./productValidation";

describe("Product Validation Tests", () => {
  const validProduct = {
    name: "Laptop Dell",
    price: 15000000,
    quantity: 10,
    description: "Mo ta ngan.",
    category: "Electronics",
  };

  // --- a) Unit Tests cho validateProduct() (3 diem) ---

  // Test product name validation
  test("TC1: Product name rong nen tra ve loi", () => {
    const product = { ...validProduct, name: "" };
    const errors = validateProduct(product);
    expect(errors.name).toBe("Ten san pham khong duoc de trong");
    expect(Object.keys(errors).length).toBe(1);
  });

  test("TC2: Product name qua ngan (< 3 ky tu) nen tra ve loi", () => {
    const product = { ...validProduct, name: "Ab" };
    const errors = validateProduct(product);
    expect(errors.name).toBe("Ten san pham phai co it nhat 3 ky tu");
  });

  // Test price validation (boundary tests)
  test("TC3: Price la so am nen tra ve loi", () => {
    const product = { ...validProduct, price: -1000 };
    const errors = validateProduct(product);
    expect(errors.price).toBe("Gia san pham phai lon hon hoac bang 0");
  });

  test("TC4: Price vuot qua Max Value (999,999,999) nen tra ve loi", () => {
    const product = { ...validProduct, price: 1000000000 };
    const errors = validateProduct(product);
    expect(errors.price).toBe("Gia san pham khong duoc vuot qua 999,999,999");
  });

  // Test quantity validation
  test("TC5: Quantity la so am nen tra ve loi", () => {
    const product = { ...validProduct, quantity: -5 };
    const errors = validateProduct(product);
    expect(errors.quantity).toBe("So luong phai lon hon hoac bang 0");
  });

  test("TC6: Quantity vuot qua Max Value (99,999) nen tra ve loi", () => {
    const product = { ...validProduct, quantity: 100000 };
    const errors = validateProduct(product);
    expect(errors.quantity).toBe("So luong khong duoc vuot qua 99,999");
  });

  // Test description length
  test("TC7: Description vuot qua 500 ky tu nen tra ve loi", () => {
    // Chuoi 501 ky tu
    const longDescription = "d".repeat(501);
    const product = { ...validProduct, description: longDescription };
    const errors = validateProduct(product);
    expect(errors.description).toBe("Mo ta khong duoc vuot qua 500 ky tu");
  });

  // Test category validation
  test("TC8: Category khong thuoc danh sach co san nen tra ve loi", () => {
    const product = { ...validProduct, category: "Invalid Category" };
    const errors = validateProduct(product);
    expect(errors.category).toBe("Danh muc khong hop le");
  });

  // Product name vuot qua 100 ky tu
  test("TC16: Product name > 100 ky tu nen tra ve loi", () => {
    const product = { ...validProduct, name: "A".repeat(101) };
    const errors = validateProduct(product);
    expect(errors.name).toBe("Ten san pham khong duoc vuot qua 100 ky tu");
  });

  // Price bi bo trong (undefined)
  test("TC17: Price undefined nen tra ve loi", () => {
    const product = { ...validProduct, price: undefined };
    const errors = validateProduct(product);
    expect(errors.price).toBe("Gia san pham khong duoc de trong");
  });

  // Quantity bi bo trong (undefined)
  test("TC18: Quantity undefined nen tra ve loi", () => {
    const product = { ...validProduct, quantity: undefined };
    const errors = validateProduct(product);
    expect(errors.quantity).toBe("So luong khong duoc de trong");
  });
});

\end{lstlisting}
\end{tcolorbox}
% ==================================================================================================================================

b) Viết tests cho Product form component (1 điểm)
% ================================================================================================================================
\begin{tcolorbox}[breakable, colback=green!5,colframe=black!50,title=ProductForm.test.js]
\begin{lstlisting}
  import { render, screen, fireEvent, waitFor } from "@testing-library/react";
import * as productService from "../services/productService";
import ProductForm from "./ProductForm";

// Mock service
jest.mock("../services/productService");

describe("ProductForm Component Tests", () => {
  beforeEach(() => {
    jest.clearAllMocks();
  });

  test("TC10: Hien thi loi validation khi submit du lieu khong hop le", async () => {
    render(<ProductForm />);
    fireEvent.change(screen.getByTestId("product-price-input"), {
      target: { value: -1000, name: "price" },
    });
    fireEvent.click(screen.getByTestId("submit-btn"));

    await waitFor(() => {
      expect(
        screen.getByText("Vui long kiem tra lai cac truong bi loi.")
      ).toBeInTheDocument();
      expect(
        screen.getByText("Gia san pham phai lon hon hoac bang 0")
      ).toBeInTheDocument();
    });
  });

  test("TC11: Submit thanh cong khi them san pham moi", async () => {
    productService.createProduct.mockResolvedValueOnce({});
    render(<ProductForm />);
    fireEvent.change(screen.getByTestId("product-name-input"), {
      target: { value: "Laptop Dell", name: "name" },
    });
    fireEvent.change(screen.getByTestId("product-price-input"), {
      target: { value: 1000, name: "price" },
    });
    fireEvent.change(screen.getByTestId("product-quantity-input"), {
      target: { value: 5, name: "quantity" },
    });
    fireEvent.change(screen.getByTestId("product-description-input"), {
      target: { value: "Mo ta ngan", name: "description" },
    });

    await waitFor(() => {
      expect(screen.getByText("Them san pham thanh cong")).toBeInTheDocument();
    });
  });

  test("TC12: Submit that bai khi API loi", async () => {
    productService.createProduct.mockRejectedValueOnce(new Error("API Error"));
    render(<ProductForm />);
    fireEvent.change(screen.getByTestId("product-name-input"), {
      target: { value: "Laptop Dell", name: "name" },
    });
    fireEvent.change(screen.getByTestId("product-price-input"), {
      target: { value: 1000, name: "price" },
    });
    fireEvent.click(screen.getByTestId("submit-btn"));

    await waitFor(() => {
      expect(screen.getByText("API Error")).toBeInTheDocument();
    });
  });

  test("TC13: Khi edit san pham, goi getProductById va hien thi du lieu", async () => {
    productService.getProductById.mockResolvedValueOnce({
      name: "Book ABC",
      price: 200,
      quantity: 2,
      description: "Sach hay",
      category: "Books",
    });

    render(<ProductForm productIdToEdit={1} />);

    await waitFor(() => {
      expect(screen.getByDisplayValue("Book ABC")).toBeInTheDocument();
      expect(screen.getByDisplayValue("200")).toBeInTheDocument();
      expect(screen.getByDisplayValue("2")).toBeInTheDocument();
      expect(screen.getByDisplayValue("Sach hay")).toBeInTheDocument();
    });
  });

  test("TC14: Khi edit san pham nhung API loi, hien thi thong bao loi", async () => {
    productService.getProductById.mockRejectedValueOnce(new Error("Not Found"));
    render(<ProductForm productIdToEdit={999} />);

\begin{lstlisting}
    await waitFor(() => {
      expect(
        screen.getByText("Loi khi tai du lieu san pham.")
      ).toBeInTheDocument();
    });
  });
});


\end{lstlisting}
\end{tcolorbox}
% ===============================================================================================================================
c) Coverage $\geq 90\%$ (1 điểm)



\subsubsection{Backend Unit Tests - Product Service}

Viết unit tests cho \texttt{ProductService}:

\textbf{Yêu cầu:}
% ==============================================================================================================================
a) Test CRUD operations:
% =============================================================================================================================
\begin{tcolorbox}[breakable, colback=green!5,colframe=black!50,title=Test createProduct()]
\begin{lstlisting}
    @Test
    @DisplayName("TC1: Tao san pham moi thanh cong")
    void testCreateProductSuccess() {
        ProductDto productDto = new ProductDto("Laptop ABC", 15000000L, 10, "Electronics");
        Product savedProduct = new Product(1L, "Laptop ABC", 15000000L, 10, "Electronics");

        when(productRepository.save(any(Product.class))).thenReturn(savedProduct);

        ProductDto result = productService.createProduct(productDto);

        assertNotNull(result);
        assertEquals(1L, result.getId());
        assertEquals("Laptop ABC", result.getName());

        verify(productRepository, times(1)).save(any(Product.class));
    }

    @Test
    @DisplayName("TC2: Lay san pham theo ID ton tai")
    void testGetProductByIdFound() {
        Long productId = 1L;
        Product mockProduct = new Product(productId, "Mouse XYZ", 200000L, 50, "Peripherals");

        when(productRepository.findById(productId)).thenReturn(Optional.of(mockProduct));

        ProductDto result = productService.getProductById(productId);

        assertNotNull(result);
        assertEquals("Mouse XYZ", result.getName());

        // VERIFY
        verify(productRepository, times(1)).findById(productId); // Bat buoc verify theo yeu cau bai tap
    }

    @Test
    @DisplayName("TC3: Lay san pham theo ID khong ton tai -> throw exception")
    void testGetProductByIdNotFound() {
        Long productId = 99L;
        when(productRepository.findById(productId)).thenReturn(Optional.empty());

        RuntimeException ex = assertThrows(RuntimeException.class,
                () -> productService.getProductById(productId));

        assertEquals("Product not found with id: " + productId, ex.getMessage());
        verify(productRepository, times(1)).findById(productId);
    }

    @Test
    @DisplayName("TC4: Cap nhat san pham thanh cong")
    void testUpdateProductSuccess() {
        Long productId = 1L;
        Product existing = new Product(productId, "Old Name", 1000L, 5, "OldCat");
        ProductDto updateDto = new ProductDto("New Name", 2000L, 10, "NewCat");
        Product updated = new Product(productId, "New Name", 2000L, 10, "NewCat");

        when(productRepository.findById(productId)).thenReturn(Optional.of(existing));
        when(productRepository.save(any(Product.class))).thenReturn(updated);

        ProductDto result = productService.updateProduct(productId, updateDto);

        assertNotNull(result);
        assertEquals("New Name", result.getName());
        assertEquals(2000L, result.getPrice());
        verify(productRepository, times(1)).findById(productId);
        verify(productRepository, times(1)).save(existing);
    }

    @Test
    @DisplayName("TC5: Xoa san pham thanh cong")
    void testDeleteProductSuccess() {
        Long productId = 1L;
        when(productRepository.existsById(productId)).thenReturn(true);
        doNothing().when(productRepository).deleteById(productId);

        productService.deleteProduct(productId);

        verify(productRepository, times(1)).existsById(productId);
        verify(productRepository, times(1)).deleteById(productId);
    }

    @Test
    @DisplayName("TC6: Xoa san pham khong ton tai -> throw exception")
    void testDeleteProductNotFound() {
        Long productId = 99L;
        when(productRepository.existsById(productId)).thenReturn(false);

        RuntimeException ex = assertThrows(RuntimeException.class,
                () -> productService.deleteProduct(productId));

        assertEquals("Product not found with id: " + productId, ex.getMessage());
        verify(productRepository, times(1)).existsById(productId);
        verify(productRepository, never()).deleteById(productId);
    }

    @Test
    @DisplayName("TC7: Lay tat ca san pham (getAll) KHONG co Pagination")
    void testGetAllProductsWithoutPagination() {
        Product p1 = new Product(1L, "Laptop", 15000000L, 10, "Electronics");
        Product p2 = new Product(2L, "Mouse", 200000L, 50, "Peripherals");

        when(productRepository.findAll()).thenReturn(java.util.List.of(p1, p2));

        java.util.List<ProductDto> result = productService.getAll();

        assertEquals(2, result.size());
        assertEquals("Laptop", result.get(0).getName());
        assertEquals("Mouse", result.get(1).getName());
        verify(productRepository, times(1)).findAll();
    }

    @Test
    @DisplayName("TC9: Lay danh sach san pham co Pagination thanh cong")
    void testGetAllProductsWithFullParams() {
        // Arrange: tao data gia lap
        Product p1 = new Product(1L, "Laptop", 15000000L, 10, "Electronics");
        Product p2 = new Product(2L, "Mouse", 200000L, 50, "Peripherals");
        java.util.List<Product> content = java.util.List.of(p1, p2);

        int page = 0;
        int size = 2;
        String nameFilter = "Lap";
        String categoryFilter = "Electronics";
        String sortBy = "price";
        String sortDir = "desc";
        long totalElements = 10L;

        Page<Product> mockPage = new PageImpl<>(
                content,
                PageRequest.of(page, size, Sort.by(sortBy).descending()),
                totalElements);

        // Mock repository method filter + pagination
        when(productRepository.findByNameContainingAndCategoryContaining(
                eq(nameFilter),
                eq(categoryFilter),
                any(Pageable.class))).thenReturn(mockPage);

        // Act
        Page<ProductDto> resultPage = productService.getAll(
                page, size, nameFilter, categoryFilter, sortBy, sortDir);

        // Assert
        assertEquals(2, resultPage.getContent().size());
        assertEquals(totalElements, resultPage.getTotalElements());
        assertEquals("Laptop", resultPage.getContent().get(0).getName());
        assertEquals("Mouse", resultPage.getContent().get(1).getName());

        verify(productRepository, times(1))
                .findByNameContainingAndCategoryContaining(eq(nameFilter), eq(categoryFilter), any(Pageable.class));
    }

}
\end{lstlisting}
\end{tcolorbox}
b) Coverage $\geq 85\%$ cho ProductService

% =======================================================================================================================================

\section{Integration Testing}

\subsection{Login - Integration Testing}

\subsubsection{Frontend Component Integration}

Test tích hợp Login component với API service:
% ====================================================================================================
a) Test rendering và user interactions
% ====================================================================================================================================
\begin{tcolorbox}[breakable, colback=green!5,colframe=black!50,title=Login Component Integration Tests]
\begin{lstlisting}
import { render, screen, fireEvent, waitFor } from "@testing-library/react";
import Login from "./Login";
import * as auth from "../services/auth";

jest.mock("../services/auth");

describe("Login Component Integration Tests", () => {
  beforeEach(() => {
    jest.clearAllMocks();
  });

  test("Nen render cac thanh phan UI co ban", () => {
    render(<Login />);

    expect(screen.getByTestId("username-input")).toBeInTheDocument();
    expect(screen.getByTestId("password-input")).toBeInTheDocument();
    expect(screen.getByTestId("login-button")).toBeInTheDocument();
  });

  test("Nen hien thi loi validation khi submit form rong", async () => {
    render(<Login />);
    const submitButton = screen.getByTestId("login-button");
    fireEvent.click(submitButton);

    await waitFor(() => {
      expect(screen.getByTestId("username-error")).toHaveTextContent(
        "Ten dang nhap khong duoc de trong"
      );
      expect(screen.getByTestId("password-error")).toHaveTextContent(
        "Mat khau khong duoc de trong"
      );
      expect(auth.login).not.toHaveBeenCalled();
    });
  });

  test("Nen goi API va hien thi thong bao thanh cong voi credentials hop le", async () => {
    auth.login.mockResolvedValue({
      success: true,
      message: "Dang nhap thanh cong",
      token: "mock-token-123",
    });

    render(<Login />);
    const usernameInput = screen.getByTestId("username-input");
    const passwordInput = screen.getByTestId("password-input");
    const submitButton = screen.getByTestId("login-button");

    fireEvent.change(usernameInput, { target: { value: "testuser" } });
    fireEvent.change(passwordInput, { target: { value: "Test123" } });
    fireEvent.click(submitButton);

    await waitFor(() => {
      expect(auth.login).toHaveBeenCalledWith("testuser", "Test123");
      expect(screen.getByTestId("login-message")).toHaveTextContent(
        "Dang nhap thanh cong"
      );
    });
  });

  test("Nen hien thi thong bao loi khi dang nhap that bai tu API", async () => {
    auth.login.mockResolvedValue({
      success: false,
      message: "Ten dang nhap hoac mat khau khong dung",
    });

    render(<Login />);
    const usernameInput = screen.getByTestId("username-input");
    const passwordInput = screen.getByTestId("password-input");
    const submitButton = screen.getByTestId("login-button");

    fireEvent.change(usernameInput, { target: { value: "testuser" } });
    fireEvent.change(passwordInput, { target: { value: "WrongPass123" } });
    fireEvent.click(submitButton);

    await waitFor(() => {
      expect(auth.login).toHaveBeenCalledTimes(1);
      expect(screen.getByTestId("login-message")).toHaveTextContent(
        "Ten dang nhap hoac mat khau khong dung"
      );
    });
  });
});
\end{lstlisting}
\end{tcolorbox}
% ======================================================================================================================================
\subsection*{Kết quả nhận được}

\begin{itemize}
    \item Tất cả 04/04 test case đều chạy thành công.
    \item Các chức năng được kiểm tra bao gồm:
    \begin{itemize}
        \item Rendering các thành phần UI cơ bản.
        \item Validation form khi submit rỗng.
        \item Gọi API mock với credentials hợp lệ và xử lý response thành công.
        \item Xử lý message lỗi khi đăng nhập thất bại từ API.
    \end{itemize}
    \item Không phát hiện lỗi về giao diện, form submission hoặc xử lý response từ API.
    \item Kết quả đạt yêu cầu đề bài:
    \begin{itemize}
        \item Rendering + interactions
        \item Form submission + API mock
        \item Xử lý message thành công + message lỗi
    \end{itemize}
\end{itemize}

b) Test form submission và API calls
\begin{tcolorbox}[breakable, colback=green!5,colframe=black!50,title=Product Form Integration Tests]
\begin{lstlisting}

import { render, screen, fireEvent, waitFor } from "@testing-library/react";
import ProductForm from "./ProductForm";
import * as productService from "../services/productService";

jest.mock("../services/productService");

describe("Product Form Integration Tests", () => {
  beforeEach(() => {
    jest.clearAllMocks();
  });

  test("Hien thi loi validation khi Price la gia tri am", async () => {
    render(<ProductForm />);

    fireEvent.change(screen.getByTestId("product-name-input"), {
      target: { value: "Test Product" },
    });
    fireEvent.change(screen.getByTestId("product-price-input"), {
      target: { value: "-1000" },
    });
    fireEvent.change(screen.getByTestId("product-quantity-input"), {
      target: { value: "5" },
    });

    fireEvent.click(screen.getByTestId("submit-btn"));

    await waitFor(() => {
      expect(productService.createProduct).not.toHaveBeenCalled();
      expect(
        screen.getByText("Gia san pham phai lon hon hoac bang 0")
      ).toBeInTheDocument();
    });
  });
});

\end{lstlisting}
\end{tcolorbox}
\subsection*{Kết quả nhận được}
\begin{itemize}
    \item Thêm sản phẩm mới thành công, API được gọi đúng, thông báo hiển thị.
    \item Validation hoạt động tốt, dữ liệu không hợp lệ không gửi API.
    \item Chỉnh sửa sản phẩm tải dữ liệu chính xác, form hoạt động ổn định.
\end{itemize}
% ============================================================================================================
c) Test error handling và success messages
% ============================================================================================================
\begin{tcolorbox}[breakable, colback=green!5,colframe=black!50,title=Product Form Success and Error Handling Tests]
\begin{lstlisting}

import { render, screen, fireEvent, waitFor } from "@testing-library/react";
import ProductForm from "./ProductForm";
import * as productService from "../services/productService";

jest.mock("../services/productService");

describe("Product Form Success and Error Handling Tests", () => {
  beforeEach(() => {
    jest.clearAllMocks();
  });

  test("Hien thi thong bao success khi tao san pham thanh cong", async () => {
    productService.createProduct.mockResolvedValue({
      id: 1,
      name: "Laptop Dell",
      price: 15000000,
      quantity: 10,
      category: "Electronics",
    });

    render(<ProductForm />);

    fireEvent.change(screen.getByTestId("product-name-input"), {
      target: { value: "Laptop Dell" },
    });
    fireEvent.change(screen.getByTestId("product-price-input"), {
      target: { value: "15000000" },
    });
    fireEvent.change(screen.getByTestId("product-quantity-input"), {
      target: { value: "10" },
    });

    fireEvent.click(screen.getByTestId("submit-btn"));

    await waitFor(() => {
      expect(screen.getByText("Them san pham thanh cong")).toBeInTheDocument();
    });
  });

  test("Hien thi thong bao loi khi server tra ve error", async () => {
    productService.createProduct.mockRejectedValue({
      message: "Server loi, khong the tao san pham",
    });

    render(<ProductForm />);

    fireEvent.change(screen.getByTestId("product-name-input"), {
      target: { value: "Laptop Dell" },
    });
    fireEvent.change(screen.getByTestId("product-price-input"), {
      target: { value: "15000000" },
    });
    fireEvent.change(screen.getByTestId("product-quantity-input"), {
      target: { value: "10" },
    });

    fireEvent.click(screen.getByTestId("submit-btn"));

    await waitFor(() => {
      expect(
        screen.getByText("Server loi, khong the tao san pham")
      ).toBeInTheDocument();
    });
  });
});
\end{lstlisting}
\end{tcolorbox}

% =============================================================================================================
\subsection*{Kết quả nhận được}
\begin{itemize}
    \item Khi tạo sản phẩm thành công: thông báo \textbf{success} hiển thị đúng, API được gọi đúng.
    \item Khi server trả về lỗi: thông báo \textbf{error} hiển thị đúng, dữ liệu không bị thêm vào.
    \item Validation và interaction với form hoạt động ổn định.
\end{itemize}

% ==============================================================================================

\subsubsection{Backend API Integration}

Test API endpoints của Login với MockMvc:

\textbf{Yêu cầu:}

a) Test POST /api/auth/login endpoint
\begin{tcolorbox}[breakable, colback=green!5,colframe=black!50,title=Login API Integration Tests]
\begin{lstlisting}
package com.flogin.controller;
import static org.springframework.test.web.servlet.request.MockMvcRequestBuilders.post;
import static org.springframework.test.web.servlet.result.MockMvcResultMatchers.status;
import static org.springframework.test.web.servlet.result.MockMvcResultMatchers.jsonPath;
import static org.mockito.Mockito.when;
import static org.mockito.ArgumentMatchers.any;
import com.fasterxml.jackson.databind.ObjectMapper;
import com.flogin.dto.LoginRequest;
import com.flogin.dto.LoginResponse;
import com.flogin.service.AuthService;
import org.junit.jupiter.api.DisplayName;
import org.junit.jupiter.api.Test;
import org.springframework.beans.factory.annotation.Autowired;
import org.springframework.boot.test.autoconfigure.web.servlet.AutoConfigureMockMvc;
import org.springframework.boot.test.autoconfigure.web.servlet.WebMvcTest;
import org.springframework.boot.test.mock.mockito.MockBean;
import org.springframework.http.MediaType;
import org.springframework.test.web.servlet.MockMvc;
@WebMvcTest(AuthController.class)
@AutoConfigureMockMvc(addFilters = false)
@DisplayName("Login API Integration Tests")
class AuthControllerIntegrationTest {


    @Autowired
    private MockMvc mockMvc;

    @Autowired
    private ObjectMapper objectMapper;

    @MockBean
    private AuthService authService;

    @Test
    @DisplayName("TC1: POST /api/auth/login - Thanh cong")
    void testLoginSuccess() throws Exception {

        LoginRequest request = new LoginRequest("testuser", "Test123");

        LoginResponse mockResponse = new LoginResponse(
                true, "Dang nhap thanh cong", "mock-token-123");

        when(authService.authenticate(any(LoginRequest.class))).thenReturn(mockResponse);

        mockMvc.perform(post("/api/auth/login")
                .contentType(MediaType.APPLICATION_JSON)
                .content(objectMapper.writeValueAsString(request)))
                .andExpect(status().isOk())
                .andExpect(jsonPath("$.success").value(true))
                .andExpect(jsonPath("$.token").exists())
                .andExpect(jsonPath("$.message").value("Dang nhap thanh cong"));
    }

    @Test
    @DisplayName("TC2: POST /api/auth/login - That bai (Sai credentials)")
    void testLoginFailure() throws Exception {

        LoginRequest request = new LoginRequest("testuser", "WrongPass");

        LoginResponse mockResponse = new LoginResponse(
                false, "Ten dang nhap hoac mat khau khong dung", null);

        when(authService.authenticate(any(LoginRequest.class))).thenReturn(mockResponse);

        mockMvc.perform(post("/api/auth/login")
                .contentType(MediaType.APPLICATION_JSON)
                .content(objectMapper.writeValueAsString(request)))
                .andExpect(status().isOk())
                .andExpect(jsonPath("$.success").value(false))
                .andExpect(jsonPath("$.token").doesNotExist())
                .andExpect(jsonPath("$.message")
                        .value("Ten dang nhap hoac mat khau khong dung"));
    }

}

\end{lstlisting}
\end{tcolorbox}
\subsection*{Kết quả nhận được}
\begin{itemize}
    \item API hoạt động đúng logic khi đăng nhập thành công
    \item Xử lý lỗi chính xác với credentials sai
    \item Đảm bảo tính ổn định trước khi tích hợp với Frontend
\end{itemize}

b) Test response structure và status codes
\begin{tcolorbox}[breakable, colback=green!5,colframe=black!50,title=Test Response Structur]
\begin{lstlisting}

@Test
void testLoginSuccess() throws Exception {
    mockMvc.perform(post("/api/auth/login")
        .contentType(MediaType.APPLICATION_JSON)
        .content(objectMapper.writeValueAsString(request)))
        .andExpect(status().isOk())                     // Kiem tra status code
        .andExpect(jsonPath("$.success").value(true))   // Kiem tra cau truc JSON
        .andExpect(jsonPath("$.token").exists())
        .andExpect(jsonPath("$.message").value("Dang nhap thanh cong"));
}
\end{lstlisting}
\end{tcolorbox}

c) Test CORS và headers
\begin{tcolorbox}[breakable, colback=green!5,colframe=black!50,title=Test CORS và Headers]
\begin{lstlisting}

@Test
@DisplayName("TC4: POST /api/auth/login - Kiem tra CORS Header")
void testCorsHeader() throws Exception {
    LoginRequest request = new LoginRequest("testuser", "Test123");

    LoginResponse mockResponse = new LoginResponse(true, "Dang nhap thanh cong", "mock-token-123");
    when(authService.authenticate(any(LoginRequest.class))).thenReturn(mockResponse);

    mockMvc.perform(post("/api/auth/login")
            .contentType(MediaType.APPLICATION_JSON)
            .content(objectMapper.writeValueAsString(request))
            .header("Origin", "http://my-frontend-domain.com")) // Gui header Origin
        .andExpect(status().isOk())
        .andExpect(header().string("Access-Control-Allow-Origin", "http://my-frontend-domain.com")); // Kiem tra header CORS
}
\end{lstlisting}
\end{tcolorbox}


\subsection{Product - Integration Testing}

\subsubsection{Frontend Component Integration}

Test tích hợp Product components:

\textbf{Yêu cầu:}

a) Test ProductList component với API
\begin{tcolorbox}[breakable, colback=green!5,colframe=black!50,title=ProductList Integration Tests]
\begin{lstlisting}

import { render, screen, waitFor } from "@testing-library/react";
import * as productService from "../services/productService";
import ProductList from "./ProductList";

// Mock toan bo module productService
jest.mock("../services/productService");

describe("ProductList Integration Tests", () => {
  beforeEach(() => {
    jest.clearAllMocks();
  });

  test("TC1: Hien thi danh sach san pham thanh cong", async () => {
    // ProductList.jsx mong doi res.content la array
    productService.getProducts.mockResolvedValueOnce({
      content: [
        { id: 1, name: "Laptop ABC", price: 1000, category: "Electronics" },
        { id: 2, name: "Mouse XYZ", price: 200, category: "Electronics" }
      ],
      totalPages: 1
    });

    render(<ProductList />);

    await waitFor(() => {
      expect(screen.getByText("Laptop ABC")).toBeInTheDocument();
      expect(screen.getByText("Mouse XYZ")).toBeInTheDocument();
    });
  });

  test("TC2: Hien thi thong bao loi khi API that bai", async () => {
    productService.getProducts.mockRejectedValueOnce(new Error("API Error"));

    render(<ProductList />);
\end{lstlisting}
\end{tcolorbox}
\begin{tcolorbox}[breakable, colback=green!5,colframe=black!50,title=ProductList Integration Tests]
\begin{lstlisting}
    await waitFor(() => {
      expect(screen.getByTestId("error-message"))
        .toHaveTextContent("Khong the tai danh sach san pham");
    });
  });

  test("TC3: Hien thi empty state khi khong co san pham", async () => {
    productService.getProducts.mockResolvedValueOnce({
      content: [],
      totalPages: 0
    });

    render(<ProductList />);

    await waitFor(() => {
      expect(
        screen.getByText("Khong tim thay san pham nao.")
      ).toBeInTheDocument();
    });
  });
});
\end{lstlisting}
\end{tcolorbox}
% =====================================================================================================
\subsection*{Kết quả nhận được}
\begin{itemize}
    \item Các sản phẩm hiển thị đầy đủ và chính xác theo API.
    \item Khi API lỗi, thông báo lỗi xuất hiện đúng.
    \item Empty state hiển thị đúng khi API trả về danh sách rỗng.
\end{itemize}

% ======================================================================================
b) Test ProductForm component (create/edit)
\begin{tcolorbox}[breakable, colback=green!5,colframe=black!50,title=Product Form Integration Test]
\begin{lstlisting}

import { render, screen, fireEvent, waitFor } from "@testing-library/react";
import ProductForm from "./ProductForm";
import * as productService from "../services/productService";

jest.mock("../services/productService");

describe("Product Form Integration Tests", () => {
  beforeEach(() => {
    jest.clearAllMocks();
  });

  test("Tao san pham moi thanh cong va hien thi thong bao success", async () => {
    productService.createProduct.mockResolvedValue({
      id: 1,
      name: "Laptop Dell",
      price: 15000000,
      quantity: 10,
      category: "Electronics",
    });

    render(<ProductForm />);

    fireEvent.change(screen.getByTestId("product-name-input"), {
      target: { value: "Laptop Dell" },
    });
    fireEvent.change(screen.getByTestId("product-price-input"), {
      target: { value: "15000000" },
    });
    fireEvent.change(screen.getByTestId("product-quantity-input"), {
      target: { value: "10" },
    });

    fireEvent.click(screen.getByTestId("submit-btn"));

    await waitFor(() => {
      expect(productService.createProduct).toHaveBeenCalledTimes(1);
      expect(screen.getByText("Them san pham thanh cong")).toBeInTheDocument();
    });
  });

  test("Hien thi loi validation khi Price la gia tri am", async () => {
    render(<ProductForm />);

    fireEvent.change(screen.getByTestId("product-name-input"), {
      target: { value: "Test Product" },
    });
    fireEvent.change(screen.getByTestId("product-price-input"), {
      target: { value: "-1000" },
    });
    fireEvent.change(screen.getByTestId("product-quantity-input"), {
      target: { value: "5" },
    });

    fireEvent.click(screen.getByTestId("submit-btn"));

    await waitFor(() => {
      expect(productService.createProduct).not.toHaveBeenCalled();
      expect(
        screen.getByText("Gia san pham phai lon hon hoac bang 0")
      ).toBeInTheDocument();
    });
  });
});
\end{lstlisting}
\end{tcolorbox}
% ========================================================================================================================================
\subsection*{Kết quả nhận được}
\begin{itemize}
    \item Thêm sản phẩm mới thành công, API gọi đúng, thông báo hiển thị.
    \item Validation hoạt động tốt, dữ liệu không hợp lệ không gửi API.
    \item Edit sản phẩm tải dữ liệu chính xác, form hoạt động ổn định.
\end{itemize}
% ====================================================================================================================================
c) Test ProductDetail component
\begin{tcolorbox}[breakable, colback=green!5,colframe=black!50,title=Product Detail Integration Tests]
\begin{lstlisting}

import { render, screen, waitFor } from "@testing-library/react";
import ProductDetail from "./ProductDetail";
import * as productService from "../services/productService";

jest.mock("../services/productService");

describe("Product Detail Integration Tests", () => {
    beforeEach(() => {
        jest.clearAllMocks();
    });

    test("TC1: Nen hien thi du lieu chi tiet san pham khi tai thanh cong", async () => {
        const mockProduct = {
            id: 1,
            name: "Sach Lap trinh",
            price: 500000,
            quantity: 5,
            category: "Books",
            description: "Tai lieu hoc React va Spring Boot"
        };

        productService.getProductById.mockResolvedValue(mockProduct);

        render(<ProductDetail productId={1} />);

        expect(screen.getByTestId("loading-message")).toBeInTheDocument();

        await waitFor(() => {
            expect(productService.getProductById).toHaveBeenCalledWith(1);
            expect(screen.getByTestId("product-detail-view")).toBeInTheDocument();
            expect(screen.getByRole('heading', { name: /Chi tiet San pham:/i })).toBeInTheDocument();
            expect(screen.getByTestId("product-price")).toHaveTextContent("500.000 VND");
            expect(screen.getByText(/So luong:/i)).toBeInTheDocument();
        });
    });

\begin{lstlisting}
    test("TC2: Nen hien thi thong bao loi khi tai chi tiet san pham that bai", async () => {
        const errorMessage = "Loi mang hoac ID khong ton tai";

        productService.getProductById.mockRejectedValue(new Error(errorMessage));

        render(<ProductDetail productId={99} />);

        await waitFor(() => {
            expect(productService.getProductById).toHaveBeenCalledWith(99);
            expect(screen.getByText(`Loi khi tai chi tiet san pham: ${errorMessage}`)).toBeInTheDocument();
        });
    });
});

\end{lstlisting}
\end{tcolorbox}
% =======================================================================================================

\subsection*{Kết quả nhận được
\begin{itemize}
    \item Chi tiết sản phẩm hiển thị chính xác khi productId hợp lệ.
    \item Khi productId không hợp lệ hoặc API lỗi, thông báo lỗi xuất hiện đúng.
    \item ProductDetail tích hợp ổn định với service, không có lỗi render.
\end{itemize}
\subsection*{Kết quả nhận được}
\begin{itemize}
    \item Các API của module Product hoạt động đúng theo yêu cầu.
    \item MockMvc + Mockito mô phỏng service thành công giúp kiểm thử controller mà không cần chạy database.
    \item Test đảm bảo 5 chức năng chính: Create – Read All – Read One – Update – Delete.
    \item Test đạt bao phủ hành vi (API behavior coverage) cao, phù hợp tiêu chuẩn kiểm thử tích hợp backend.
\end{itemize}


% ===================================================================================================================
\subsubsection{API Integration}

Test các API endpoints của Product:

% ====================================================================================================================
a) Test POST /api/products (Create)
% ====================================================================================================================
\begin{tcolorbox}[breakable, colback=green!5,colframe=black!50,title=Create Product Tests ]
\begin{lstlisting}
@Test
@DisplayName("TC1: POST /api/products - Tao san pham moi thanh cong")
void testCreateProduct() throws Exception {
    ProductDto requestDto = new ProductDto("Laptop XYZ", 20000000L, 5, "Electronics");

    ProductDto responseDto = new ProductDto(1L, "Laptop XYZ", 20000000L, 5, "Electronics");

    when(productService.createProduct(any(ProductDto.class))).thenReturn(responseDto);

    mockMvc.perform(post("/api/products")
            .contentType(MediaType.APPLICATION_JSON)
            .content(objectMapper.writeValueAsString(requestDto)))

            .andExpect(status().isOk())
            .andExpect(jsonPath("$.id").exists())
            .andExpect(jsonPath("$.name").value("Laptop XYZ"))
            .andExpect(jsonPath("$.price").value(20000000.0));

    verify(productService, times(1)).createProduct(any(ProductDto.class));
}
\end{lstlisting}
\end{tcolorbox}
% =====================================================================================================================
b) Test GET /api/products (Read all)
% ============================================================================================================================
\begin{tcolorbox}[breakable, colback=green!5,colframe=black!50,title=Read All Product Tests ]
\begin{lstlisting}

@Test
@DisplayName("TC2: GET /api/products/all - Lay danh sach san pham thanh cong")
void testGetAllProducts() throws Exception {

    List<ProductDto> products = Arrays.asList(
            new ProductDto(1L, "Laptop", 15000000L, 10, "Electronics"),
            new ProductDto(2L, "Mouse", 200000L, 50, "Electronics"));

    when(productService.getAll()).thenReturn(products);

    mockMvc.perform(get("/api/products/all"))
            .andExpect(status().isOk())
            .andExpect(jsonPath("$", hasSize(2)))
            .andExpect(jsonPath("$[0].name").value("Laptop"))
            .andExpect(jsonPath("$[1].quantity").value(50));
}

\end{lstlisting}
\end{tcolorbox}
% =======================================================================================================================
c) Test GET /api/products/\{id\} (Read one)
% ==========================================================================================================================
\begin{tcolorbox}[breakable, colback=green!5,colframe=black!50,title=Read One Product Tests ]
\begin{lstlisting}
@Test
@DisplayName("TC3: GET /api/products/{id} - Lay chi tiet san pham thanh cong")
void testGetProductById() throws Exception {
    Long productId = 1L;
    ProductDto mockProduct = new ProductDto(productId, "Keyboard", 1000000L, 20, "Electronics");

    when(productService.getProductById(productId)).thenReturn(mockProduct);

    mockMvc.perform(get("/api/products/{id}", productId))
            .andExpect(status().isOk())
            .andExpect(jsonPath("$.id").value(productId))
            .andExpect(jsonPath("$.name").value("Keyboard"))
            .andExpect(jsonPath("$.price").value(1000000.0));

    verify(productService, times(1)).getProductById(productId);
}
\end{lstlisting}
\end{tcolorbox}
% ==================================================================================================================================
d) Test PUT /api/products/\{id\} (Update)
% ===========================================================================================================================
\begin{tcolorbox}[breakable, colback=green!5,colframe=black!50,title=Update Product Tests ]
\begin{lstlisting}
@Test
@DisplayName("TC4: PUT /api/products/{id} - Cap nhat san pham thanh cong")
void testUpdateProduct() throws Exception {
    Long productId = 1L;
    ProductDto updateRequest = new ProductDto("Updated Laptop", 18000000L, 8, "Electronics");
    ProductDto updatedProduct = new ProductDto(productId, "Updated Laptop", 18000000L, 8, "Electronics");

    when(productService.updateProduct(eq(productId), any(ProductDto.class))).thenReturn(updatedProduct);

    mockMvc.perform(put("/api/products/{id}", productId)
            .contentType(MediaType.APPLICATION_JSON)
            .content(objectMapper.writeValueAsString(updateRequest)))

            .andExpect(status().isOk())
            .andExpect(jsonPath("$.id").value(productId))
            .andExpect(jsonPath("$.name").value("Updated Laptop"))
            .andExpect(jsonPath("$.price").value(18000000.0));

    verify(productService, times(1)).updateProduct(eq(productId), any(ProductDto.class));
}
\end{lstlisting}
\end{tcolorbox}
% ======================================================================================================================
e) Test DELETE /api/products/\{id\} (Delete)
% ==================================================================================================================================
\begin{tcolorbox}[breakable, colback=green!5,colframe=black!50,title=Delete Product Tests ]
\begin{lstlisting}

@Test
@DisplayName("TC5: DELETE /api/products/{id} - Xoa san pham thanh cong")
void testDeleteProduct() throws Exception {
    Long productId = 1L;

    doNothing().when(productService).deleteProduct(productId);

    mockMvc.perform(delete("/api/products/{id}", productId))
            .andExpect(status().isNoContent());

    verify(productService, times(1)).deleteProduct(productId);
}

\end{lstlisting}
\end{tcolorbox}
% ==========================================================================================
\section{ Mock Testing}

\subsection{Login - Mock Testing}

\subsubsection{Frontend Mocking }

Mock external dependencies cho Login component:

\textbf{Yêu cầu:}
% ===================================================================================================================
a) Mock \texttt{authService.loginUser()}
\begin{tcolorbox}[breakable, colback=green!5,colframe=black!50,title=Mock authService.loginUser() ]
\begin{lstlisting}

import Login from "./Login";

import { validateUsername, validatePassword } from "../utils/validation";
import * as authService from "../services/auth";

jest.mock("../utils/validation");
jest.mock("../services/auth");

const mockedValidateUsername = jest.mocked(validateUsername);
const mockedValidatePassword = jest.mocked(validatePassword);
const mockedLogin = jest.mocked(authService.login);
\end{lstlisting}
\end{tcolorbox}
b) Test với mocked successful/failed responses (1 điểm)
\begin{tcolorbox}[breakable, colback=green!5,colframe=black!50,title=Test Với Mocked Successful/Failed Responses ]
\begin{lstlisting}
test("Mock: Dang nhap thanh cong", async () => {
  mockedValidateUsername.mockReturnValue("");
  mockedValidatePassword.mockReturnValue("");

  mockedLogin.mockResolvedValue({
    success: true,
    token: "Mock-token-123",
    message: "Dang nhap thanh cong",
  });

  render(<Login />);

  fireEvent.change(screen.getByTestId("username-input"), {
    target: { value: "validuser" },
  });
  fireEvent.change(screen.getByTestId("password-input"), {
    target: { value: "validpass" },
  });
});

fireEvent.click(screen.getByTestId("login-button"));

await waitFor(() => {
  expect(mockedLogin).toHaveBeenCalledTimes(1);

  expect(mockedLogin).toHaveBeenCalledWith("validuser", "validpass");

  expect(screen.getByTestId("login-message")).toHaveTextContent(
    "Dang nhap thanh cong"
  );
});
});

test("Mock: hien thi thong bao that bai tu server khi success la false", async () => {
  mockedValidateUsername.mockReturnValue("");
  mockedValidatePassword.mockReturnValue("");

  mockedLogin.mockResolvedValue({
    success: false,
    message: "Sai mat khau.",
  });

  render(<Login />);
  fireEvent.click(screen.getByTestId("login-button"));

  await waitFor(() => {
    expect(screen.getByTestId("login-message")).toHaveTextContent(
      "Sai mat khau."
    );
  });
});

\end{lstlisting}
\end{tcolorbox}
% ================================================================================================
c) Verify mock calls
% =============================================================================================
\begin{tcolorbox}[breakable, colback=green!5,colframe=black!50,title=Verify Mock Calls ]
\begin{lstlisting}
// Verify mock calls
expect(mockedLogin).toHaveBeenCalledTimes(1);
expect(mockedLogin).toHaveBeenCalledWith("validuser", "validpass");
\end{lstlisting}
\end{tcolorbox}
% ====================================================================================
\subsubsection{Backend Mocking}

Mock dependencies trong Backend tests:


% ===================================================================================
a) Mock \texttt{AuthService} với \texttt{@MockBean}
% ==================================================================================================
\begin{tcolorbox}[breakable, colback=green!5,colframe=black!50,title=Mock AuthService Với @MockBean ]
\begin{lstlisting}
@WebMvcTest(AuthController.class)
@AutoConfigureMockMvc(addFilters = false)
class AuthControllerMockTest {

@Autowired
private MockMvc mockMvc;

@MockBean
private AuthService authService;
}
\end{lstlisting}
\end{tcolorbox}


b) Test controller với mocked service (1 điểm)
\begin{tcolorbox}[breakable, colback=green!5,colframe=black!50,title=Test Controller Với Mocked Service ]
\begin{lstlisting}
@Test
@DisplayName("Mock: Controller login thanh cong voi mocked service")
void testLoginWithMockedService() throws Exception {

    LoginResponse mockResponse = new LoginResponse(
            true,
            "Success",
            "mock-token"
    );

    when(authService.authenticate(any()))
            .thenReturn(mockResponse);

    mockMvc.perform(post("/api/auth/login")
                    .contentType(MediaType.APPLICATION_JSON)
                    .content(" {\" username \":\" test \" ,\" password \":\" Pass123 \"}"))

            .andExpect(status().isOk());

    verify(authService, times(1)).authenticate(any(LoginRequest.class));
}

@Test
@DisplayName("Mock: Controller login that bai")
void testLoginFailureWithMockedService() throws Exception {
    LoginResponse mockResponse = new LoginResponse(false, "Invalid credentials", null);

    when(authService.authenticate(any(LoginRequest.class)))
            .thenReturn(mockResponse);

    mockMvc.perform(post("/api/auth/login")
                    .contentType(MediaType.APPLICATION_JSON)
                    .content(" {\" username \":\" test \" ,\" password \":\" Pass123 \"}"))

            .andExpect(status().isOk())
            .andExpect(jsonPath("$.success").value(false));

    verify(authService, times(1)).authenticate(any(LoginRequest.class));
}

\end{lstlisting}
\end{tcolorbox}
c) Verify mock interactions
\begin{tcolorbox}[breakable, colback=green!5,colframe=black!50,title=Verify Mock Interactions ]
\begin{lstlisting}
verify(authService, times(1)).authenticate(any(LoginRequest.class));
\end{lstlisting}
\end{tcolorbox}

\subsection{Product - Mock Testing}

\subsubsection{Frontend Mocking}

Mock \texttt{ProductService} trong component tests:


a) Mock CRUD operations
\begin{tcolorbox}[breakable, colback=green!5,colframe=black!50,title= Mock CRUD Operations ]
\begin{lstlisting}
import ProductList from "./ProductList";
import ProductForm from "./ProductForm";

import * as productService from "../services/productService";
import * as productValidation from "../utils/productValidation";

jest.mock("../services/productService");
jest.mock("../utils/productValidation");

const mockedGetProducts = jest.mocked(productService.getProducts);
const mockedGetProductById = jest.mocked(productService.getProductById);
const mockedCreateProduct = jest.mocked(productService.createProduct);
const mockedUpdateProduct = jest.mocked(productService.updateProduct);
const mockedDeleteProduct = jest.mocked(productService.deleteProduct);
const mockedValidateProduct = jest.mocked(productValidation.validateProduct);

\end{lstlisting}
\end{tcolorbox}
b) Test success và failure scenarios
\begin{tcolorbox}[breakable, colback=green!5,colframe=black!50,title=Test Success Và Failure Scenarios ]
\begin{lstlisting}
describe("Product Mock Tests (CRUD)", () => {
  beforeEach(() => {
    jest.clearAllMocks();
    mockedValidateProduct.mockReturnValue({});
  });

  const mockList = [
    { id: 1, name: "Laptop Dell", price: 15000000 },
    { id: 2, name: "Chuot Logitech", price: 500000 },
  ];

  test("CREATE: Them san pham moi thanh cong", async () => {
    mockedCreateProduct.mockResolvedValue({ success: true });
    render(<ProductForm />);

    fireEvent.change(screen.getByTestId("product-name-input"), {
      target: { value: "iPhone 15" },
    });
    fireEvent.change(screen.getByTestId("product-price-input"), {
      target: { value: "30000000" },
    });
    fireEvent.change(screen.getByTestId("product-quantity-input"), {
      target: { value: "10" },
    });
    fireEvent.change(screen.getByTestId("product-category-select"), {
      target: { value: "Electronics" },
    });

    fireEvent.click(screen.getByTestId("submit-btn"));

    await waitFor(() => {
      expect(mockedCreateProduct).toHaveBeenCalledWith(
        expect.objectContaining({
          name: "iPhone 15",
          price: 30000000,
        })
      );
      expect(screen.getByText("Them san pham thanh cong")).toBeInTheDocument();
    });
  });

  test("READ: Tai danh sach that bai", async () => {
    mockedGetProducts.mockRejectedValue(new Error("Network Error"));
    render(<ProductList />);

    await waitFor(() => {
      expect(mockedGetProducts).toHaveBeenCalledTimes(1);
      expect(screen.getByTestId("error-message")).toHaveTextContent(
        "Khong the tai danh sach san pham"
      );
    });
  });
});

\end{lstlisting}
\end{tcolorbox}

c) Verify all mock calls
\begin{tcolorbox}[breakable, colback=green!5,colframe=black!50,title= Verify All Mock Calls ]
\begin{lstlisting}
// CREATE
expect(mockedCreateProduct).toHaveBeenCalledWith(
  expect.objectContaining({
    name: "iPhone 15",
    price: 30000000,
  })
);

// READ
expect(mockedGetProducts).toHaveBeenCalledTimes(1);

// DELETE
expect(mockedDeleteProduct).toHaveBeenCalledWith(2);

// UPDATE
expect(mockedGetProductById).toHaveBeenCalledWith(1);
expect(mockedUpdateProduct).toHaveBeenCalledWith(
  1,
  expect.objectContaining({
    name: "Laptop Dell",
    price: 14000000,
  })
);
\end{lstlisting}
\end{tcolorbox}



\subsubsection{Backend Mocking}

Mock \texttt{ProductRepository} trong service tests:


a) Mock \texttt{ProductRepository}
\begin{tcolorbox}[breakable, colback=green!5,colframe=black!50,title=Mock ProductRepository ]
\begin{lstlisting}
@ExtendWith(MockitoExtension.class)
class ProductServiceMockTest {
    @Mock
    private ProductRepository productRepository;

    @InjectMocks
    private ProductServiceImpl productService;
}

\end{lstlisting}
\end{tcolorbox}
b) Test service layer với mocked repository
\begin{tcolorbox}[breakable, colback=green!5,colframe=black!50,title=Test Service Layer Với Mocked Repository ]
\begin{lstlisting}
@Test
@DisplayName("Mock: Lay toan bo danh sach san pham")
void testGetAll() {
    Product p1 = new Product(1L, "A", 100L, 1, "Cat1");
    Product p2 = new Product(2L, "B", 200L, 2, "Cat2");

    when(productRepository.findAll()).thenReturn(java.util.Arrays.asList(p1, p2));

    java.util.List<ProductDto> result = productService.getAll();

    assertEquals(2, result.size());
    assertEquals("A", result.get(0).getName());

    verify(productRepository, times(1)).findAll();
}

@Test
@DisplayName("Mock: Tao san pham moi thanh cong")
void testCreateProduct() {
    ProductDto inputDto = new ProductDto("Mouse", 200000L, 50, "Accessories");

    Product savedEntity = new Product();
    savedEntity.setId(2L);
    savedEntity.setName("Mouse");
    savedEntity.setPrice(200000L);

    when(productRepository.save(any(Product.class))).thenReturn(savedEntity);

    ProductDto result = productService.createProduct(inputDto);

    assertNotNull(result);
    assertEquals(2L, result.getId());
    assertEquals("Mouse", result.getName());

    verify(productRepository, times(1)).save(any(Product.class));
}
\end{lstlisting}
\end{tcolorbox}
c) Verify repository interactions
\begin{tcolorbox}[breakable, colback=green!5,colframe=black!50,title= Verify Repository Interactions ]
\begin{lstlisting}
// GET ALL
verify(productRepository, times(1)).findAll();

// CREATE
verify(productRepository, times(1)).save(any(Product.class));


\end{lstlisting}
\end{tcolorbox}
% ========================================================================================
\section{Automation Testing và CI/CD}

\subsection{ Login - E2E Automation Testing}

\subsubsection{Setup và Configuration}


\begin{itemize}
    \item Cài đặt Cypress: \textbf{npm install cypress --save-dev}

% \begin{mdframed}[backgroundcolor=green!10,hidealllines=true,
%   innerleftmargin=10pt,innerrightmargin=10pt,
%   innertopmargin=10pt,innerbottommargin=10pt]

% \begin{lstlisting}[caption={Cài đặt Cypress}, captionpos=b]


% \end{lstlisting}
% \end{mdframed}

    \item Cấu hình test environment
\begin{tcolorbox}[breakable, colback=green!5,colframe=black!50,title=Test Environment ]
\begin{lstlisting}
const { defineConfig } = require("cypress");

module.exports = defineConfig({
  e2e: {
    baseUrl: "http://localhost:5173",
    setupNodeEvents(on, config) {
      // implement node event listeners here
    },
  },
});

\end{lstlisting}
\end{tcolorbox}


    \item Setup Page Object Model

% \begin{itemize}

    % ---------- ITEM CON 1 ----------
    % \item LoginPage.js

\begin{tcolorbox}[breakable, colback=green!5,colframe=black!50,title=LoginPage.js]
\begin{lstlisting}
class LoginPage {
  // Selectors
  usernameInput = '[data-testid="username-input"]';
  passwordInput = '[data-testid="password-input"]';
  loginButton = '[data-testid="login-button"]';
  loginMessage = '[data-testid="login-message"]';
  usernameError = '[data-testid="username-error"]';

  // Actions
  visit() {
    cy.visit("/login");
  }

  fillLoginForm(username, password) {
    cy.get(this.usernameInput).type(username);
    cy.get(this.passwordInput).type(password);
  }

  submit() {
    cy.get(this.loginButton).click();
  }

  getLoginMessage() {
    return cy.get(this.loginMessage);
  }

  getUsernameError() {
    return cy.get(this.usernameError);
  }
}

export default new LoginPage();

\end{lstlisting}
\end{tcolorbox}


    % ---------- ITEM CON 2 ----------
%     \item ProductPage.js

% \begin{tcolorbox}[breakable, colback=green!5,colframe=black!50,title=ProductPage.js ]
% \begin{lstlisting}
% class ProductPage {
%   visit() {
%     cy.visit("/product-list");
%   }

%   clickAddNew() {
%     cy.get('[data-testid="add-new-btn"]').click();
%   }

%   openFormModal() {
%     cy.get('[data-testid="modal-overlay"]').should(
%       "be.visible"
%     );
%   }

%   closeFormModal() {
%     cy.get('[data-testid="close-modal-btn"]').click();
%   }

%   fillProductForm(product) {
%     cy.get('[data-testid="product-name-input"]').clear().type(product.name);
%     cy.get('[data-testid="product-price-input"]').clear().type(product.price);
%   }
% \end{lstlisting}
% \end{tcolorbox}
% \begin{tcolorbox}[breakable, colback=green!5,colframe=black!50,title=ProductPage.js ]
% \begin{lstlisting}
%   submitForm() {
%     cy.get('[data-testid="submit-btn"]').click();
%   }

%   getSuccessMessage() {
%     return cy.get('[data-testid="success-message"]');
%   }

%   getProductInList(name) {
%     return cy.contains('[data-testid="product-item"]', name);
%   }

%   clickEditButton(name) {
%     this.getProductInList(name).find('[data-testid^="edit-btn-"]').click();
%   }

%   clickDeleteButton(name) {
%     this.getProductInList(name).find('[data-testid^="delete-btn-"]').click();
%   }

%   fillSearchInput(keyword) {
%     cy.get('[data-testid="search-input"]').clear().type(keyword);
%   }

%   selectSortByNewest() {
%     cy.get('[data-testid="sort-select"]').select("id-desc");
%   }
% }

% export default ProductPage;

% \end{lstlisting}
% \end{tcolorbox}

% \end{itemize}


\end{itemize}


\subsubsection{E2E Test Scenarios cho Login}

% Viết automated tests cho toàn bộ login flow:

% a) Test complete login flow
\begin{tcolorbox}[breakable, colback=green!5,colframe=black!50,title=Automated tests cho toàn bộ login flow]
\begin{lstlisting}
import LoginPage from "../pages/LoginPage";

describe("Login E2E Tests (5.1.2)", () => {
  beforeEach(() => {
    LoginPage.visit();
  });

  it("TC1: Hien thi form login va kiem tra tuong tac co ban", () => {
    cy.get(LoginPage.usernameInput)
      .should("be.visible")
      .and("have.attr", "type", "text");
    cy.get(LoginPage.passwordInput)
      .should("be.visible")
      .and("have.attr", "type", "password");
    cy.get(LoginPage.loginButton).should("be.visible").and("be.enabled");

    cy.get(LoginPage.usernameInput).type("temp").should("have.value", "temp");
    cy.get(LoginPage.usernameInput).clear().should("have.value", "");
  });

  it("TC2: Dang nhap thanh cong voi credentials hop le", () => {
    LoginPage.fillLoginForm("admin", "Test123");
    LoginPage.submit();

    LoginPage.getLoginMessage().should("contain", "Dang nhap thanh cong");
  });

  it("TC3: Dang nhap that bai voi credentials sai (API Error)", () => {
    LoginPage.fillLoginForm("testuser", "WrongPass456");
    LoginPage.submit();

    LoginPage.getLoginMessage().should("contain", "Sai thong tin dang nhap");
  });

  it("TC4: Hien thi loi validation khi Username qua ngan (Boundary Test)", () => {
    LoginPage.fillLoginForm("ab", "Test123");
    LoginPage.submit();

    LoginPage.getUsernameError()
      .should("be.visible")
      .and("contain", "phai co it nhat 3 ky tu");
  });

  it("TC5: Hien thi loi validation khi Password khong co chu hoac so (Negative Test)", () => {
    LoginPage.fillLoginForm("testuser", "abcdef");
    LoginPage.submit();

    cy.get('[data-testid="password-error"]')
      .should("be.visible")
      .and("contain", "Mat khau phai chua ca chu va so");
  });

  it("TC6: Ngan chan ky tu dac biet khong hop le trong Username (Edge Case)", () => {
    LoginPage.fillLoginForm("user!@#", "Test123");
    LoginPage.submit();

    LoginPage.getUsernameError()
      .should("be.visible")
      .and("contain", "Ten dang nhap chi duoc chua a-z, A-Z, 0-9");
  });
});
\end{lstlisting}
\end{tcolorbox}
% \begin{tcolorbox}[breakable, colback=green!5,colframe=black!50,title=Delete Product Tests]
% \begin{lstlisting}


% \end{lstlisting}
% \end{tcolorbox}

% b) Test validation messages (0.5 điểm)
% \begin{mdframed}[backgroundcolor=green!10,hidealllines=true,
%   innerleftmargin=10pt,innerrightmargin=10pt,
%   innertopmargin=10pt,innerbottommargin=10pt]

% \begin{lstlisting}[caption={Các test validation messages}, captionpos=b]

%         VIET TEST VAO
%         VIET TEST VAO
%         VIET TEST VAO

% \end{lstlisting}
% \end{mdframed}
% c) Test success/error flows (0.5 điểm)
% \begin{mdframed}[backgroundcolor=green!10,hidealllines=true,
%   innerleftmargin=10pt,innerrightmargin=10pt,
%   innertopmargin=10pt,innerbottommargin=10pt]

% \begin{lstlisting}[caption={Các test success/error flows}, captionpos=b]

%         VIET TEST VAO
%         VIET TEST VAO
%         VIET TEST VAO

% \end{lstlisting}
% \end{mdframed}

% d) Test UI elements interactions (0.5 điểm)
% \begin{mdframed}[backgroundcolor=green!10,hidealllines=true,
%   innerleftmargin=10pt,innerrightmargin=10pt,
%   innertopmargin=10pt,innerbottommargin=10pt]

% \begin{lstlisting}[caption={Các test UI elements interactions }, captionpos=b]

%         VIET TEST VAO
%         VIET TEST VAO
%         VIET TEST VAO

% \end{lstlisting}
% \end{mdframed}

\subsubsection{CI/CD Integration cho Login Tests}


\begin{itemize}
    \item Tạo GitHub Actions workflow
\begin{tcolorbox}[breakable, colback=green!5,colframe=black!50,title=Tạo GitHub Actions Workflow ]
\begin{lstlisting}
name: Login Tests CI

on:
  push:
    branches: [main, develop]
  pull_request:
    branches: [main]

jobs:
  test-login:
    runs-on: ubuntu-latest

    steps:
      - uses: actions/checkout@v2

      - name: Setup Java
        uses: actions/setup-java@v2
        with:
          java-version: "17"
          distribution: "temurin"

      - name: Make Maven wrapper executable
        run: chmod +x ./backend/mvnw

      - name: Start Backend in Background (Port 8080)
        run: |
          cd backend
          ./mvnw spring-boot:run &
          sleep 20

      - name: Setup Node.js
        uses: actions/setup-node@v2
        with:
          node-version: "18"

      - name: Install dependencies
        run: |
          cd frontend
          npm install

      - name: Run Login Unit Tests
        run: |
          cd frontend
          npm test -- --coverage --testPathPatterns Login
        - name: Start Frontend in Background (Port 5173)
        run: |
          cd frontend
          npm run build
          npm install -g serve
          serve -s dist -l 5173 &
          sleep 10

      - name: Run Login E2E Tests
        run: |
          cd frontend
          npm run test:e2e:spec -- "cypress/e2e/login.cy.js"
        env:
          BASE_URL: http://localhost:5173

      - name: Upload Test Reports
        uses: actions/upload-artifact@v4
        # uses: actions/upload-artifact@v3
        with:
          name: test-reports
          path: frontend/reports/
\end{lstlisting}
\end{tcolorbox}
% \begin{tcolorbox}[breakable, colback=green!5,colframe=black!50,title=Tạo GitHub Actions Workflow ]
% \begin{lstlisting}


% \end{lstlisting}
% \end{tcolorbox}
    \item Generate test reports
\begin{tcolorbox}[breakable, colback=green!5,colframe=black!50,title=Generate test reports]
\begin{lstlisting}
<?xml version="1.0" encoding="UTF-8"?>
<testsuites name="Mocha Tests" time="4.347" tests="6" failures="0">
  <testsuite name="Root Suite" timestamp="2025-11-30T11:15:21" tests="0" file="cypress\e2e\login.cy.js" time="0.000" failures="0">
  </testsuite>
  <testsuite name="Login E2E Tests (5.1.2)" timestamp="2025-11-30T11:15:21" tests="6" time="4.343" failures="0">
    <testcase name="Login E2E Tests (5.1.2) TC1: Hiển thị form login và kiểm tra tương tác cơ bản" time="0.617" classname="TC1: Hiển thị form login và kiểm tra tương tác cơ bản">
    </testcase>
    <testcase name="Login E2E Tests (5.1.2) TC2: Đăng nhập thành công với credentials hợp lệ" time="0.988" classname="TC2: Đăng nhập thành công với credentials hợp lệ">
    </testcase>
    <testcase name="Login E2E Tests (5.1.2) TC3: Đăng nhập thất bại với credentials sai (API Error)" time="0.739" classname="TC3: Đăng nhập thất bại với credentials sai (API Error)">
    </testcase>
    <testcase name="Login E2E Tests (5.1.2) TC4: Hiển thị lỗi validation khi Username quá ngắn (Boundary Test)" time="0.584" classname="TC4: Hiển thị lỗi validation khi Username quá ngắn (Boundary Test)">
    </testcase>
    <testcase name="Login E2E Tests (5.1.2) TC5: Hiển thị lỗi validation khi Password không có chữ hoặc số (Negative Test)" time="0.630" classname="TC5: Hiển thị lỗi validation khi Password không có chữ hoặc số (Negative Test)">
    </testcase>
    <testcase name="Login E2E Tests (5.1.2) TC6: Ngăn chặn ký tự đặc biệt không hợp lệ trong Username (Edge Case)" time="0.643" classname="TC6: Ngăn chặn ký tự đặc biệt không hợp lệ trong Username (Edge Case)">
    </testcase>
  </testsuite>
</testsuites>
\end{lstlisting}
\end{tcolorbox}
%     \item Generate test reports
% \begin{mdframed}[backgroundcolor=green!10,hidealllines=true,
%   innerleftmargin=10pt,innerrightmargin=10pt,
%   innertopmargin=10pt,innerbottommargin=10pt]

% \begin{lstlisting}[caption={Generate test reports }, captionpos=b]

%         VIET TEST VAO
%         VIET TEST VAO
%         VIET TEST VAO

% \end{lstlisting}
% \end{mdframed}
\end{itemize}



\subsection{Product - E2E Automation Testing}

\subsubsection{Setup Page Object Model}

Implement POM cho Product pages:
\begin{tcolorbox}[breakable, colback=green!5,colframe=black!50,title=Implement POM Cho Product Pages ]
\begin{lstlisting}
class ProductPage {
  visit() {
    cy.visit("/product-list");
  }

  clickAddNew() {
    cy.get('[data-testid="add-new-btn"]').click();
  }

  openFormModal() {
    cy.get('[data-testid="modal-overlay"]').should(
      "be.visible"
    );
  }

  closeFormModal() {
    cy.get('[data-testid="close-modal-btn"]').click();
  }

  fillProductForm(product) {
    cy.get('[data-testid="product-name-input"]').clear().type(product.name);
    cy.get('[data-testid="product-price-input"]').clear().type(product.price);
  }

  submitForm() {
    cy.get('[data-testid="submit-btn"]').click();
  }

  getSuccessMessage() {
    return cy.get('[data-testid="success-message"]');
  }

  getProductInList(name) {
    return cy.contains('[data-testid="product-item"]', name);
  }

  clickEditButton(name) {
    this.getProductInList(name).find('[data-testid^="edit-btn-"]').click();
  }

    clickDeleteButton(name) {
    this.getProductInList(name).find('[data-testid^="delete-btn-"]').click();
  }

  fillSearchInput(keyword) {
    cy.get('[data-testid="search-input"]').clear().type(keyword);
  }

  selectSortByNewest() {
    cy.get('[data-testid="sort-select"]').select("id-desc");
  }
}

export default ProductPage;

\end{lstlisting}
\end{tcolorbox}
% \begin{tcolorbox}[breakable, colback=green!5,colframe=black!50,title=Implement POM Cho Product Pages ]
% \begin{lstlisting}


% \end{lstlisting}
% \end{tcolorbox}

\subsubsection{E2E Test Scenarios cho Product}

Viết automated tests cho CRUD operations:
\begin{tcolorbox}[breakable, colback=green!5,colframe=black!50,title=Automated Tests ]
\begin{lstlisting}
// Dua tren Listing 16
import ProductPage from "../pages/ProductPage";

const NEW_PRODUCT_DATA = {
  name: "Laptop HP Spectre X360 (Test)",
  price: "35000000",
  quantity: "3",
};

const UPDATE_PRODUCT_DATA = {
  price: "33000000",
};

describe("Product E2E Tests (Su dung data-testid)", () => {
  const productPage = new ProductPage();

  beforeEach(() => {
    productPage.visit();
  });

  // a) Test Create product flow (0.5 diem)
  it("Nen tao san pham moi thanh cong (Create)", () => {
    productPage.clickAddNew();
    productPage.openFormModal();

    productPage.fillProductForm(NEW_PRODUCT_DATA);
    productPage.submitForm();

    productPage.getSuccessMessage().should("contain", "Them moi thanh cong!");
    productPage.getProductInList(NEW_PRODUCT_DATA.name).should("exist");
  });

  // b) Test Read/List products (0.5 diem)
  it("Nen hien thi san pham vua tao trong danh sach", () => {
    // Dam bao san pham ton tai truoc khi test
    productPage.getProductInList(NEW_PRODUCT_DATA.name).should("be.visible");
  });

  // e) Test Search/Filter functionality (0.5 diem)
  it("Nen tim kiem san pham theo tu khoa", () => {
    productPage.fillSearchInput("HP Spectre");

    productPage.getProductInList("Laptop HP Spectre X360").should("exist");
    productPage.getProductInList("Chuot Logitech").should("not.exist");
  });

  // c) Test Update product (0.5 diem)
  it("Nen cap nhat san pham thanh cong (Update)", () => {
    // 1. Click nut Sua
    productPage.clickEditButton(NEW_PRODUCT_DATA.name);
    productPage.openFormModal();

    // 2. Cap nhat gia
    productPage.fillProductForm({
      name: NEW_PRODUCT_DATA.name,
      price: UPDATE_PRODUCT_DATA.price,
    });

    // 3. Submit form
    productPage.submitForm();

    // 4. Kiem tra thong bao thanh cong
    productPage.getSuccessMessage().should("contain", "Cap nhat thanh cong!");

    // 5. Kiem tra gia tri moi trong danh sach
    const formattedPrice = UPDATE_PRODUCT_DATA.price.toLocaleString();
    productPage
      .getProductInList(NEW_PRODUCT_DATA.name)
      .invoke("text")
      .should("contain", UPDATE_PRODUCT_DATA.price.toLocaleString());
  });

  // d) Test Delete product (0.5 diem)
  it("Nen xoa san pham thanh cong (Delete)", () => {
    productPage.clickDeleteButton(NEW_PRODUCT_DATA.name);
    productPage.getSuccessMessage().should("contain", "Xoa thanh cong!");
    productPage.getProductInList(NEW_PRODUCT_DATA.name).should("not.exist");
  });
});

\end{lstlisting}
\end{tcolorbox}

\subsubsection{CI/CD Integration}

\begin{itemize}
    \item Tạo GitHub Actions workflow
\begin{tcolorbox}[breakable, colback=green!5,colframe=black!50,title=Tạo GitHub Actions Workflow ]
\begin{lstlisting}
name: Product Tests CI
on:
  push:
    branches: [main, develop]
  pull_request:
    branches: [main]

jobs:
  test-product:
    runs-on: ubuntu-latest

    steps:
      - uses: actions/checkout@v2

      # --- Setup Java & Backend ---
      - name: Setup Java
        uses: actions/setup-java@v2
        with:
          java-version: "17"
          distribution: "temurin"

      - name: Make Maven wrapper executable
        run: chmod +x ./backend/mvnw

      - name: Start Backend in Background (Port 8080)
        run: |
          cd backend
          ./mvnw spring-boot:run &
          sleep 20

      # --- Setup Node & Frontend ---
      - name: Setup Node.js
        uses: actions/setup-node@v2
        with:
          node-version: "18"

      - name: Install dependencies
        run: |
          cd frontend
          npm install

      - name: Start Frontend in Background (Port 5173)
        run: |
          cd frontend
          npm run build
          npm install -g serve
          serve -s dist -l 5173 &
          sleep 10

      # --- Run Product E2E Tests ---
      - name: Run Product E2E Tests
        run: |
          cd frontend
          npm run test:e2e:spec -- "cypress/e2e/product.cy.js"
        env:
          BASE_URL: http://localhost:5173

      # --- Upload Product Test Reports ---
      - name: Upload Product Test Reports
        # uses: actions/upload-artifact@v3
        uses: actions/upload-artifact@v4
        with:
          name: product-test-reports
          path: frontend/reports/

\end{lstlisting}
\end{tcolorbox}
    \item Generate test reports
\begin{tcolorbox}[breakable, colback=green!5,colframe=black!50,title=Generate test reports]
\begin{lstlisting}
<?xml version="1.0" encoding="UTF-8"?>
<testsuites name="Mocha Tests" time="4.937" tests="5" failures="0">
  <testsuite name="Root Suite" timestamp="2025-11-30T11:15:27" tests="0" file="cypress\e2e\product.cy.js" time="0.000" failures="0">
  </testsuite>
  <testsuite name="Product E2E Tests (Su dung data-testid)" timestamp="2025-11-30T11:15:27" tests="5" time="4.933" failures="0">
    <testcase name="Product E2E Tests (Su dung data-testid) Nen tao san pham moi thanh cong (Create)" time="1.423" classname="Nen tao san pham moi thanh cong (Create)">
    </testcase>
    <testcase name="Product E2E Tests (Su dung data-testid) Nen hien thi san pham vua tao trong danh sach" time="1.547" classname="Nen hien thi san pham vua tao trong danh sach">
    </testcase>
    <testcase name="Product E2E Tests (Su dung data-testid) Nen tim kiem san pham theo tu khoa" time="0.461" classname="Nen tim kiem san pham theo tu khoa">
    </testcase>
    <testcase name="Product E2E Tests (Su dung data-testid) Nen cap nhat san pham thanh cong (Update)" time="1.202" classname="Nen cap nhat san pham thanh cong (Update)">
    </testcase>
    <testcase name="Product E2E Tests (Su dung data-testid) Nen xoa san pham thanh cong (Delete)" time="0.190" classname="Nen xoa san pham thanh cong (Delete)">
    </testcase>
  </testsuite>
</testsuites>

\end{lstlisting}
\end{tcolorbox}
\end{itemize}
\end{enumerate}
\newpage
\section{Phần Mở Rộng}

\subsection{Performance Testing}

% ================== SETUP ======================
\subsubsection{Setup k6 cho performance testing}

\paragraph{Công cụ sử dụng:}
\begin{itemize}
    \item \textbf{k6 (v0.55.0)} – công cụ kiểm thử tải mã nguồn mở, viết bằng Go, sử dụng JavaScript cho scripting.
    \item \textbf{Lý do lựa chọn:} k6 tiêu thụ tài nguyên thấp, hỗ trợ tạo lượng lớn Virtual Users (VUs), dễ dàng tích hợp vào CI/CD.
\end{itemize}

\paragraph{Cài đặt công cụ:}
Cài đặt trên Windows bằng Winget:
\begin{verbatim}
winget install k6 --source winget
\end{verbatim}
\begin{figure}[h]
    \centering
    \includegraphics[width=0.8\textwidth]{2.png}
    \caption{Minh hoạ quá trình cài đặt k6}
\end{figure}

\begin{verbatim}
k6 version
\end{verbatim}
\begin{figure}[h]
    \centering
    \includegraphics[width=0.8\textwidth]{4.png}
    \caption{Minh hoạ quá trình cài đặt k6}
\end{figure}



\paragraph{Môi trường kiểm thử:}
\begin{itemize}
    \item \textbf{Hardware:} Máy trạm Windows SSD.
    \item \textbf{Backend:} Spring Boot chạy Tomcat (Port 8080).
    \item \textbf{Database:} H2 In-memory.
    \item \textbf{Network:} Localhost – loại bỏ nhiễu mạng.
\end{itemize}

\paragraph{Chuẩn bị kịch bản:}
\begin{verbatim}
cd performance test
\end{verbatim}

\begin{figure}[h]
    \centering
    \includegraphics[width=0.8\textwidth]{3.png}
    \caption{Thư mục chứa các file test}
\end{figure}

% ================== LOGIN TEST ======================
\subsubsection{Performance tests cho Login API}

\paragraph{Load Test}
Kịch bản trong file \texttt{login-test.js} gồm:
\begin{itemize}
    \item 0 → 100 users trong 30 giây.
    \item 100 → 500 users trong 1 phút.
\end{itemize}

\paragraph{Stress Test}
\begin{itemize}
    \item Đẩy tải lên 1000 concurrent users trong 1 phút.
\end{itemize}

\paragraph{Response Time Analysis}
\begin{itemize}
    \item \textbf{http\_req\_duration}: mục tiêu $<200ms$.
    \item \textbf{checks}: Xác minh đăng nhập thành công (Status 200 + Token).
\end{itemize}

\paragraph{Thực thi:}
\begin{verbatim}
k6 run login-test.js
\end{verbatim}

\begin{figure}[h]
    \centering
    \includegraphics[width=0.8\textwidth]{5.png}
    \caption{Kết quả test Login API}
\end{figure}

% ================== PRODUCT TEST ======================
\subsubsection{Performance tests cho Product API}

\paragraph{Load Scenario}
\begin{itemize}
    \item Giả lập 500 users gọi liên tục API \texttt{/api/products} trong 1 phút.
\end{itemize}

\paragraph{Thực thi:}
\begin{verbatim}
k6 run product-test.js
\end{verbatim}

\begin{figure}[H]
    \centering
    \includegraphics[width=0.8\textwidth]{6.png}
    \caption{Kết quả test Product API}
\end{figure}

% ================== RECOMMENDATIONS ======================
\subsubsection{Phân tích kết quả và recommendations}

\paragraph{Phân tích:}
\begin{itemize}
    \item 500 users: Hệ thống phản hồi ổn định.
    \item 1000 users: Thời gian phản hồi tăng nhưng server không crash.
\end{itemize}

\paragraph{Recommendations:}
\begin{itemize}
    \item Thêm \textbf{Redis Cache} cho API sản phẩm.
    \item Cấu hình \textbf{Connection Pool} khi deploy Production.
    \item Nếu mở rộng, cân nhắc \textbf{Load Balancing}.
\end{itemize}

\end{document}
